% content/c03-hilbertraum/de.tex — Kapitel 3: Werkzeugkasten Hilbertraum
% (deutsche Leitfassung)
%
% Stand: der Stoff der zweiten und dritten Vorlesung (latex/lectures/de/vl02.tex, „Vorlesung:
% Hilbertraum“, und latex/lectures/de/vl03.tex, „Vorlesung: Hilbertraum, Basis, Operatoren“, ohne
% deren Unterabschnitt Operatoren).  Das Kapitel wächst später um vl04 (Operatoren, mit dem
% Unterabschnitt Operatoren von vl03) und vl05 (Spur, Diagonalisierung).  Neue Abschnitte kommen
% als \section c03.s8, c03.s9, ... zwischen den letzten Kern-Abschnitt und die Fundament-Blöcke
% am Dateiende; die Fundament-Blöcke bleiben am Ende.
%
% Wohin jede Zeile von vl02 und vl03 gegangen ist: evidence/migrations/maps/vl02-map.json und
%   vl03-map.json (geprüft mit evidence/migrations/check_vl02_map.py und check_vl03_map.py).
% Fehler in vl02: evidence/migrations/M-0013.json bis M-0016.json; in vl03: M-0017.json bis
%   M-0024.json.
% Definitionen: Registry D75 bis D80 (vl02), D81 bis D86 und D14 (vl03).  Herleitungen:
%   evidence/derivations/c03/DERIVATION.md (Abschnitt „Lecture 3“ für vl03), Anker = Claim-ID;
%   Zahlen nachgerechnet von evidence/derivations/c03/check_c03.py.  Seitenbild-Zitate aus
%   Reed-Simon I: evidence/provenance.json P-0140 bis P-0150 (vl02, Prüfquittung
%   evidence/receipts/reviews/R-0140.json) und P-0152 bis P-0164 (vl03, Prüfquittung
%   evidence/receipts/reviews/R-0152.json).
% Prüfungen der dritten Vorlesung: docs/review/lectures/REVIEW-lec3.md (Stimme und Richtigkeit)
%   und evidence/derivations/c03/REVIEW-lec3.md (Seitenbilder, Mathematik, Migrationsakten);
%   alle Befunde beider Tabellen sind in einer Reparaturrunde (2026-10-21) geschlossen, die
%   Stellen sind im Text mit der Zeile der Prüfung kommentiert.
% Stimme: Sätze aus vl02.tex und vl03.tex sind beibehalten, wo sie stimmen; Bemerkungen aus den
%   Videos (transcripts/clean/v02.md ab 76:43, v03.md, v04.md bis 46:45; Zeitmarken geprüft an
%   transcripts/raw/) sind mit Zeitmarke kommentiert.
%
% D4-Brücke (PLAN.md 1.3 und 14.7): Der Kern ist endlichdimensional.  C^n und das Skalarprodukt
% auf C^n sind Kern, ebenso die allgemeinen Axiome des Skalarprodukts, Norm, Schwarz- und
% Dreiecksungleichung, D14 und Basen mit endlicher Indexmenge; l^2 und L^2(R) stehen in den
% Vertiefungen c03.v1 und c03.v2, Basen mit unendlicher Indexmenge in c03.v4, Vollständigkeit,
% Äquivalenzklassen und Cauchy-Schwarz in C^n und l^2 im Fundament.
% GAP-TJO: Der Wortlaut der Brücke in Abschnitt c03.s3 ist ein Vorschlag des Kapitelautors
% (G-0033 (b); die Brücke in C5 ist G-0004).  Offene Punkte der dritten Vorlesung: G-0072.

\chapter{Werkzeugkasten Hilbertraum}\label{c03}

% Stimme: vl02.tex:4 und v03 [00:24] („das ist der mathematische Teil der Quantenmechanik“).
Um die Zustände der Quantenmechanik zu beschreiben, wird der \emph{Hilbertraum} eine große
Rolle spielen. Er ist der mathematische Kern der Quantenmechanik: Ohne ihn können wir über
ihre Objekte nicht einmal reden. Um diesen Raum zu definieren, benötigen wir erst
\emph{komplexe Vektorräume}. Solche Räume werden ganz analog zum reellen Fall definiert, den
Sie aus der Theoretischen Physik~A kennen \srcref{tpc-script-v0.2}{vl02, Einleitung}.

% Stimme: vl02.tex:6, v02 [77:16]-[77:49].
Warum hier \emph{komplexe} und nicht reelle Vektoren auftreten, wissen wir nicht: Die Natur
ist nun einmal so gebaut, wir müssen damit leben, und eine bessere Alternative wurde bisher
nicht gefunden. Diese Tatsache hat sehr wichtige physikalische Folgen, die Sie im Lauf des
Kurses sehen werden.

In Kapitel~\ref{c02} ist das Qubit schon vorgekommen, mit seiner Zustandskugel. Wie die
Zustände der Quantenmechanik genau aus einem Hilbertraum gebaut werden, sagt Kapitel~4. Dieses
Kapitel stellt das Werkzeug bereit, und zwar so, dass man es in die Hand nehmen kann: mit
Spalten aus komplexen Zahlen.

\section{Komplexe Vektoren}\label{c03.s1}

% Stimme: v02 [79:05]-[83:31].
Glücklicherweise ist die Definition einfach. Komplexe Zahlen kennen Sie, Spaltenvektoren
auch; jetzt werden beide zusammengebracht.

\tpcdef{D75}

Ein komplexer Vektor ist also nichts anderes als eine Liste komplexer Zahlen. An der Tafel
schreibe ich ihn fett und mit einem Strich darunter; gedruckt ist er nur fett, $\vb{v}$.
Multipliziert man einen komplexen Vektor mit einem Skalar oder addiert man zwei, so kommt
wieder ein komplexer Vektor heraus.

% Stimme: vl02.tex:55, v02 [85:05].
Für den Vektorraum selbst genügen diese beiden Rechenregeln. Ein Skalarprodukt braucht man
dafür nicht: Kreuz-Operation und Skalarprodukt aus dem nächsten Abschnitt sind
Zusatzstruktur, und erst mit ihnen wird aus $\mathbb{C}^n$ ein Hilbertraum.

% Stimme: v03 [08:52]-[10:31].
Der einfachste Fall, $\mathbb{C}^1$, ist langweilig: Das sind einfach die komplexen Zahlen. Das
erste interessante Beispiel ist $\mathbb{C}^2$. Dieser Raum ist rechnerisch ganz einfach und
physikalisch trotzdem zentral: Zweidimensionale komplexe Vektoren beschreiben die Zustände
eines Qubits, des Grundbausteins eines Quantencomputers, und man braucht sie, um den Spin
eines Elektrons zu beschreiben \srcref{tpc-script-v0.2}{vl17, Abschnitt Spin}. Genauer
beschreibt ein Vektor einen \emph{reinen} Zustand; gemischte Zustände, wie die Punkte im
Inneren der Zustandskugel aus Kapitel~\ref{c02}, brauchen Dichtematrizen, und die kommen in
Kapitel~4 \srcref{tpc-script-v0.2}{vl06, Postulat 2}.

\begin{example}[Rechnen in $\mathbb{C}^2$]\label{c03.ex.c2}\claim{c03.cl.c2}{derived}
  Der Vektorraum $\mathbb{C}^2$ besteht aus allen Vektoren mit zwei komplexen Einträgen. Seien
  \begin{equation}\label{c03.eq.c2-vektoren}
    \vb{v} = \begin{pmatrix} 1+i \\ 1-i \end{pmatrix},
    \qquad
    \vb{w} = \begin{pmatrix} 1 \\ -1 \end{pmatrix}
  \end{equation}
  zwei Vektoren; der erste hat ausdrücklich komplexe Einträge, um zu betonen, dass das
  erlaubt ist. Eintragsweise gerechnet ist
  \begin{equation}\label{c03.eq.c2-rechnung}
    \vb{v} + \vb{w} = \begin{pmatrix} 2+i \\ -i \end{pmatrix},
    \qquad
    i\,\vb{v} = \begin{pmatrix} i + i^2 \\ i - i^2 \end{pmatrix}
              = \begin{pmatrix} -1+i \\ 1+i \end{pmatrix} .
  \end{equation}
\end{example}

\section{Kreuz-Operation und Skalarprodukt}\label{c03.s2}

% Stimme: v03 [02:34]-[07:47].
Ein komplexer Vektorraum allein trägt noch keine weitere Struktur. Um zum Hilbertraum zu
kommen, brauchen wir zusätzliche Operationen. Die erste ist die \emph{Kreuz-Operation}: Aus
einem Spaltenvektor wird ein Zeilenvektor, und dabei werden die Einträge komplex
konjugiert. Mit ihr lässt sich sofort das Skalarprodukt einführen: Man wendet die
Kreuz-Operation auf den ersten Vektor an und multipliziert die Zeile mit der Spalte des
zweiten. Heraus kommt immer eine komplexe Zahl.

\tpcdef{D76}

\begin{tpcfigure}{c03.fig.kreuz}
  \demo{c03.demo.kreuz}
  \caption{Die Kreuz-Operation und das Skalarprodukt. (a) Die Einträge eines Vektors in
    $\mathbb{C}^2$ als Punkte der komplexen Ebene, hier $v_1 = 2+i$ und $v_2 = -1+2i$
    (gefüllt). Die Kreuz-Operation spiegelt jeden Eintrag an der reellen Achse (offene
    Kreise, $v_1^* = 2-i$ und $v_2^* = -1-2i$) und legt die Spalte als Zeile hin. (b) Das
    Skalarprodukt aus Beispiel~\ref{c03.ex.skalarprodukt}: Die Zeile $\vb{v}\adj$ wird
    Eintrag für Eintrag mit der Spalte $\vb{w}$ multipliziert, und die Produkte werden
    addiert.}
\end{tpcfigure}

% Stimme: v03 [04:40] („das muss man nicht, aber es ist hilfreich später“).
Die komplexe Konjugation in der Kreuz-Operation muss man nicht machen; man könnte eine Spalte
auch einfach umlegen. Sie zahlt sich aber aus, sobald man einen Vektor mit sich selbst
multipliziert. Weil $v_j^* v_j = \abs{v_j}^2 = (\operatorname{Re} v_j)^2 + (\operatorname{Im} v_j)^2$
ist, gilt
\begin{equation}\label{c03.eq.selbst}\claim{c03.cl.selbst}{derived}
  \vb{v}\adj\vb{v} = \sum_{j=1}^{n} v_j^* v_j = \sum_{j=1}^{n} \abs{v_j}^2 \;\geq\; 0 .
\end{equation}
Das Skalarprodukt eines Vektors mit sich selbst ist also eine reelle Zahl und nie negativ.

\begin{example}[Ohne Konjugation]\label{c03.ex.ohne-stern}\claim{c03.cl.ohne-stern}{derived}
  Sei $\vb{v} = \begin{pmatrix} 1 \\ i \end{pmatrix}$. Ohne Konjugation ergäbe Zeile mal
  Spalte $1 \cdot 1 + i \cdot i = 1 + i^2 = 0$, obwohl $\vb{v}$ nicht der Nullvektor ist. Mit
  der Kreuz-Operation ist
  \begin{equation}\label{c03.eq.ohne-stern}
    \vb{v}\adj\vb{v} = 1 \cdot 1 + (-i) \cdot i = 1 + 1 = 2 .
  \end{equation}
  Ohne den Stern hätte ein Vektor, der nicht null ist, die „Länge“ null.
\end{example}

% Stimme: v03 [07:16]-[07:47]; vl03.tex:132.
Das Skalarprodukt ist die Zusatzstruktur, mit der wir später in der Quantenmechanik
\emph{Wahrscheinlichkeiten} berechnen. Die Quantenmechanik ist eine probabilistische Theorie
(Kapitel~\ref{c01}) und braucht deshalb eine Methode, Wahrscheinlichkeiten auszurechnen; das
Skalarprodukt wird diese Methode sein \srcref{tpc-script-v0.2}{vl03, Abschnitt Skalarprodukte}.

\begin{example}[Ein Skalarprodukt in $\mathbb{C}^2$]\label{c03.ex.skalarprodukt}\claim{c03.cl.vw}{derived}
  Für die Vektoren aus Gleichung~\eqref{c03.eq.c2-vektoren} rechnet man in zwei Schritten,
  erst die Kreuz-Operation, dann Zeile mal Spalte:
  \begin{equation}\label{c03.eq.vw}
    \vb{v}\adj\vb{w}
    = \begin{pmatrix} 1-i & 1+i \end{pmatrix} \begin{pmatrix} 1 \\ -1 \end{pmatrix}
    = (1-i) - (1+i) = -2i .
  \end{equation}
  Ist es egal, auf welchen der beiden Vektoren man das Kreuz anwendet? Nein: Der
  \emph{erste} Vektor bekommt immer das Kreuz. In der anderen Reihenfolge ist die
  Konjugation hier trivial, weil $\vb{w}$ reell ist, und es kommt etwas anderes heraus:
  \begin{equation}\label{c03.eq.wv}
    \vb{w}\adj\vb{v}
    = \begin{pmatrix} 1 & -1 \end{pmatrix} \begin{pmatrix} 1+i \\ 1-i \end{pmatrix}
    = (1+i) - (1-i) = 2i .
  \end{equation}
  Die beiden Ergebnisse sind zueinander komplex konjugiert. Die Reihenfolge ist also wichtig,
  aber der Unterschied ist kontrolliert.
\end{example}

Dass das kein Zufall ist, liest man an der Summe ab. Die folgenden Regeln gelten für jedes
$n$.

\begin{proposition}[Rechenregeln des Skalarprodukts in $\mathbb{C}^n$]\label{c03.prop.skalarprodukt}\claim[proof-included]{c03.cl.skalarprodukt}{derived}
  Seien $\vb{v}, \vb{w}, \vb{u} \in \mathbb{C}^n$ und $\alpha, \beta \in \mathbb{C}$. Dann gilt:
  \begin{enumerate}
    \item Vertauschen konjugiert: $\vb{w}\adj\vb{v} = (\vb{v}\adj\vb{w})^*$.
    \item Linear im zweiten Faktor:
      $\vb{v}\adj(\alpha\vb{w} + \beta\vb{u}) = \alpha\, \vb{v}\adj\vb{w} + \beta\, \vb{v}\adj\vb{u}$.
    \item Antilinear im ersten Faktor:
      $(\alpha\vb{v} + \beta\vb{u})\adj\vb{w} = \alpha^*\, \vb{v}\adj\vb{w} + \beta^*\, \vb{u}\adj\vb{w}$.
    \item Positivität: $\vb{v}\adj\vb{v} \geq 0$, und $\vb{v}\adj\vb{v} = 0$ gilt nur für
      $\vb{v} = 0$.
  \end{enumerate}
\end{proposition}
\begin{proof}
  Zu~1: Die Konjugierte einer Summe ist die Summe der Konjugierten, und $(z^*)^* = z$, also
  \begin{equation}\label{c03.eq.vertauschen}
    (\vb{v}\adj\vb{w})^* = \sum_{j=1}^{n} (v_j^* w_j)^* = \sum_{j=1}^{n} v_j\, w_j^*
    = \sum_{j=1}^{n} w_j^* v_j = \vb{w}\adj\vb{v} .
  \end{equation}
  Zu~2: $\sum_j v_j^* (\alpha w_j + \beta u_j) = \alpha \sum_j v_j^* w_j + \beta \sum_j v_j^* u_j$.
  Zu~3: Der $j$-te Eintrag von $(\alpha\vb{v} + \beta\vb{u})\adj$ ist
  $(\alpha v_j + \beta u_j)^* = \alpha^* v_j^* + \beta^* u_j^*$; weiter wie bei~2. Zu~4:
  Gleichung~\eqref{c03.eq.selbst}. Eine Summe von Zahlen $\abs{v_j}^2 \geq 0$ ist nur dann
  null, wenn jeder Summand null ist, also jedes $v_j = 0$.
\end{proof}

Beim Herausziehen eines Faktors aus dem ersten Vektor kommt also ein Stern dazu. Das ist die
Konvention dieses Kurses; manche Bücher machen es umgekehrt und nennen das Skalarprodukt
linear im ersten Vektor \srcref{reed-simon-1}{S.~37}.

\begin{exercise}[Rechnen mit dem Kreuz]\label{c03.exc.rechnen}\claim{c03.cl.rechnen}{derived}
  Seien $\vb{v} = \begin{pmatrix} i \\ 2 \end{pmatrix}$ und
  $\vb{w} = \begin{pmatrix} 1 \\ i \end{pmatrix}$. Berechnen Sie $3\vb{v}$,
  $\vb{v} + \vb{w}$, $\vb{v}\adj\vb{w}$, $\vb{w}\adj\vb{v}$ und $\vb{v}\adj\vb{v}$, und prüfen
  Sie Regel~1 aus Proposition~\ref{c03.prop.skalarprodukt}.
\end{exercise}
\begin{solution}
  Es ist $3\vb{v} = \begin{pmatrix} 3i \\ 6 \end{pmatrix}$ und
  $\vb{v} + \vb{w} = \begin{pmatrix} 1+i \\ 2+i \end{pmatrix}$. Mit
  $\vb{v}\adj = \begin{pmatrix} -i & 2 \end{pmatrix}$ und
  $\vb{w}\adj = \begin{pmatrix} 1 & -i \end{pmatrix}$ ist
  \begin{equation}\label{c03.eq.rechnen}
    \vb{v}\adj\vb{w} = -i + 2i = i,
    \qquad
    \vb{w}\adj\vb{v} = i - 2i = -i,
    \qquad
    \vb{v}\adj\vb{v} = (-i)\,i + 4 = 5 .
  \end{equation}
  Also ist $\vb{w}\adj\vb{v} = -i = (\vb{v}\adj\vb{w})^*$, wie Regel~1 sagt.
\end{solution}

\section{Unendlich viele Einträge}\label{c03.s3}

% Stimme: v03 [16:54]-[18:28]; vl02.tex:85.
% GAP-TJO: In der Vorlesung (v03 [17:26]-[18:28]) folgt hier der Satz, es sei offen und
% umstritten, ob man das Universum mit endlich vielen Freiheitsgraden modellieren könne.  Er
% hat keine Quelle auf der Platte und steht deshalb nicht im Text (G-0033 (a)).
An dieser Stelle könnte man endlichdimensionale Hilberträume schon definieren. Ich verschiebe
das aber bewusst, weil unendlichdimensionale Vektorräume in der Physik eine sehr wichtige
Rolle spielen. Der Grund ist physikalisch: Unsere besten Modelle sind Felder, und Felder
haben unendlich viele Freiheitsgrade, wie Sie aus der Theoretischen Physik~B wissen
\srcref{tpc-script-v0.2}{vl02, Absatz über unendlich dimensionale Vektorräume}.

% GAP-TJO: Brückenabsatz des Kapitelautors (G-0033 (b)).
Der Kern dieses Kurses bleibt trotzdem bis einschließlich Kapitel~4 endlichdimensional:
Alles, was dort über Zustände und Messungen gesagt wird, spielt in $\mathbb{C}^n$. Die
unendlichdimensionalen Räume betreten den Kern erst in Kapitel~5, beim Teilchen auf der
Geraden. Wer schon jetzt sehen will, was schiefgeht, wenn man in den Formeln einfach
$n = \infty$ setzt, und wie man es repariert, findet das in zwei Vertiefungen. Die erste
handelt von Spalten mit unendlich vielen Einträgen\zoomref{c03.v1}.

\begin{vertiefung}[id=c03.v1,
    titel={Spalten mit unendlich vielen Einträgen},
    wozu={Setzt man in den Formeln von Abschnitt~\ref{c03.s2} einfach $n = \infty$, so kann
      das Skalarprodukt divergieren. Hier wird gezeigt, wie man das repariert: Man lässt
      nur Vektoren zu, deren Skalarprodukt mit sich selbst endlich ist. Diese Vektoren
      bilden den Raum $\elltwo$.},
    voraussetzung={Abschnitt~\ref{c03.s2}; konvergente Reihen. Der Kern bis Kapitel~4
      braucht den Block nicht.},
    zurueck={Ein Vektor mit unendlich vielen Einträgen ist zugelassen, wenn die Summe der
      Betragsquadrate seiner Einträge endlich ist. Für solche Vektoren ist das
      Skalarprodukt als unendliche Summe erklärt und hat dieselben Rechenregeln wie in
      $\mathbb{C}^n$.}]
  % Stimme: vl02.tex:85-96; v03 [18:59]-[21:50].
  Man kann sich unendlichdimensionale Vektoren als Spalten mit unendlich vielen Einträgen
  vorstellen, also als unendliche Listen komplexer Zahlen:
  \begin{equation*}
    \vb{v} = \begin{pmatrix} v_1 \\ v_2 \\ v_3 \\ \vdots \end{pmatrix} .
  \end{equation*}
  Addition und Multiplikation mit Skalaren gehen Eintrag für Eintrag wie in $\mathbb{C}^n$.
  Für das Skalarprodukt ist die naheliegende Verallgemeinerung die Reihe
  $\vb{v}\adj\vb{w} = \sum_{j=1}^{\infty} v_j^* w_j$: In der Regel setzt man in allen
  Gleichungen $n = \infty$. Hier können aber Probleme auftreten, Konvergenzprobleme. Eines
  zeigt sich sofort.

  \begin{example}[Der Vektor mit lauter Einsen]\label{c03.ex.einsen}\claim{c03.cl.einsen}{derived}
    Der Vektor mit lauter Einsen, $v_j = 1$ für alle $j$, ist als unendliche Liste
    komplexer Zahlen völlig legal. Sein Skalarprodukt mit sich selbst divergiert aber: Die
    Teilsummen sind
    \begin{equation}\label{c03.eq.einsen}
      \sum_{j=1}^{N} \abs{v_j}^2 = \sum_{j=1}^{N} 1 = N \;\longrightarrow\; \infty
      \qquad (N \to \infty) .
    \end{equation}
  \end{example}

  % Stimme: v03 [22:54]-[24:28] („in der mathematischen Literatur sehr kompliziert
  % beschrieben, aber physikalisch gesehen ist es einfach“).
  Wenn wir die Zustände physikalischer Systeme mit solchen Vektoren beschreiben wollen, ist
  das ein Problem, spätestens bei den Wahrscheinlichkeiten. Die Antwort ist denkbar
  einfach: Wir lassen solche Vektoren nicht zu. In der mathematischen Literatur wird das
  sehr kompliziert beschrieben; physikalisch gesehen ist es einfach.

  \tpcdef{D77}

  % Stimme: v03 [26:41]-[27:47].
  Die Bezeichnung $\elltwo$ werden wir später kaum brauchen, Sie sollen sie aber einmal
  gesehen haben: Die $2$ kommt vom Quadrat in der Summe, und summiert wird über alle
  natürlichen Zahlen. Reed und Simon schreiben den Raum mit einer tiefgestellten $2$
  \srcref{reed-simon-1}{S.~40}. Dass die zugelassenen Vektoren wieder einen Vektorraum
  bilden, ist nicht offensichtlich, aber wahr:

  \begin{proposition}[$\elltwo$ ist ein Vektorraum mit Skalarprodukt]\label{c03.prop.elltwo}\claim[proof-included]{c03.cl.elltwo}{derived}
    Seien $\vb{v}, \vb{w} \in \elltwo$ und $a \in \mathbb{C}$. Dann gilt:
    \begin{enumerate}
      \item $a\vb{v} \in \elltwo$ und $\vb{v} + \vb{w} \in \elltwo$.
      \item Die Reihe $\vb{v}\adj\vb{w} = \sum_{j=1}^{\infty} v_j^* w_j$ konvergiert absolut,
        und $\sum_{j=1}^{\infty} \abs{v_j^* w_j} \leq \tfrac12\, \vb{v}\adj\vb{v} + \tfrac12\, \vb{w}\adj\vb{w}$.
      \item Die vier Regeln von Proposition~\ref{c03.prop.skalarprodukt} gelten in $\elltwo$.
    \end{enumerate}
  \end{proposition}
  \begin{proof}
    Für komplexe Zahlen $a, b$ ist $0 \leq (\abs{a} - \abs{b})^2$ und
    $\abs{a + b}^2 = \abs{a}^2 + \abs{b}^2 + 2 \operatorname{Re}(a^* b)$ mit
    $\operatorname{Re}(a^* b) \leq \abs{a}\,\abs{b}$. Daraus folgen
    \begin{equation}\label{c03.eq.ab}
      \abs{a^* b} = \abs{a}\,\abs{b} \leq \tfrac12 \abs{a}^2 + \tfrac12 \abs{b}^2,
      \qquad
      \abs{a + b}^2 \leq 2\abs{a}^2 + 2\abs{b}^2 .
    \end{equation}
    Zu~1: $\sum_{j} \abs{a v_j}^2 = \abs{a}^2 \sum_j \abs{v_j}^2$. Nach der zweiten
    Ungleichung ist jede Teilsumme $\sum_{j \leq N} \abs{v_j + w_j}^2$ höchstens
    $2\, \vb{v}\adj\vb{v} + 2\, \vb{w}\adj\vb{w}$. Die Teilsummen wachsen mit $N$ und sind
    beschränkt, also konvergieren sie. Zu~2: genauso mit der ersten Ungleichung; eine absolut
    konvergente Reihe komplexer Zahlen konvergiert. Zu~3: Die Regeln gelten für jede
    Teilsumme bis $N$, denn das ist eine Rechnung in $\mathbb{C}^N$. Grenzwerte vertragen
    sich mit Summen, Vielfachen und Konjugation. Ist $\vb{v}\adj\vb{v} = 0$, so ist jede
    Teilsumme $\sum_{j \leq N} \abs{v_j}^2$ null, also jedes $v_j = 0$.
  \end{proof}

  % Stimme: vl02.tex:112; v03 [26:41].
  Mit dieser Bedingung kann bei der Berechnung von Wahrscheinlichkeiten nichts mehr
  unendlich groß werden. Eine schärfere Schranke als die in Teil~2, die
  Cauchy-Schwarz-Ungleichung, steht im Fundament\zoomref{c03.f1}.

  Dass die Einträge gegen null gehen, reicht übrigens nicht. Es kommt darauf an, wie schnell.

  \begin{example}[Es kommt auf die Rate an]\label{c03.ex.rate}\claim{c03.cl.rate}{derived}
    Der Vektor mit $v_j = 1/j$ liegt in $\elltwo$, der Vektor mit $u_j = 1/\sqrt{j}$ nicht,
    obwohl auch seine Einträge gegen null gehen. Für jedes $N \geq 2$ ist nämlich
    \begin{equation}\label{c03.eq.rate}
      \begin{aligned}
        \sum_{j=1}^{N} \frac{1}{j^2}
          &\leq 1 + \sum_{j=2}^{N} \Bigl( \frac{1}{j-1} - \frac{1}{j} \Bigr)
           = 2 - \frac{1}{N} < 2, \\
        \sum_{j=1}^{N} \abs{u_j}^2 = \sum_{j=1}^{N} \frac{1}{j}
          &\geq \int_{1}^{N+1} \frac{\dd x}{x} = \ln(N+1) \;\longrightarrow\; \infty .
      \end{aligned}
    \end{equation}
    In der ersten Zeile ist $1/j^2 \leq 1/((j-1)j) = 1/(j-1) - 1/j$, und die Summe rechts
    schiebt sich zusammen. In der zweiten liegt der Summand $1/j$ über der Fläche unter
    $1/x$ zwischen $j$ und $j+1$, weil $1/x$ fällt. Die Teilsummen der Betragsquadrate von
    $\vb{u}$ wachsen also über alle Grenzen, mindestens wie ein Logarithmus.
  \end{example}

  \begin{tpcfigure}{c03.fig.partialsummen}
    \demo{c03.demo.partialsummen}
    \caption{Die Zulassungsprüfung für $\elltwo$: die Teilsummen
      $\sum_{j=1}^{N} \abs{v_j}^2$ für $N = 1, \dots, 25$, auf vier Stellen gerundet; die
      Kurve für lauter Einsen ist nur bis $N = 5$ gezeichnet. Für
      $v_j = 1/j$ (Punkte, durchgezogen) bleiben sie unter der Schranke $2$ aus
      Gleichung~\eqref{c03.eq.rate} (dünne Linie): Der Vektor ist zugelassen. Für
      $v_j = 1/\sqrt{j}$ (Punkte, gestrichelt) wachsen sie über alle Grenzen, mindestens wie
      $\ln(N+1)$, für den
      Vektor mit lauter Einsen (Punkte, strichpunktiert) wie $N$: Beide sind nicht
      zugelassen, obwohl beim mittleren die Einträge gegen null gehen. In $\mathbb{C}^n$ hat
      die Summe nur $n$ Glieder, und es gibt nichts zu prüfen.}
  \end{tpcfigure}

  % Stimme: v03 [28:28]-[29:29]; vl02.tex:139.
  % GAP-TJO: In der Vorlesung (v03 [29:58]) und im Skript (vl02.tex:139) folgt das Beispiel
  % „ein Photon in einer Glasfaser“; was damit gemeint ist, entscheidet TJO (G-0033 (e)).
  Dieser Raum ist nicht nur Mathematik. $\elltwo$ ist der Zustandsraum des harmonischen
  Oszillators, den wir in Kapitel~7 genau ansehen. Als Quantenmechaniker arbeite ich
  praktisch jeden Tag mit diesen beiden Vektorräumen, mit $\mathbb{C}^2$ und mit $\elltwo$.
\end{vertiefung}

% Stimme: v03 [30:23]-[36:22].
Der dritte Raum, den wir brauchen, besteht nicht aus Spalten, sondern aus Funktionen: Jeder
Stelle $x$ der reellen Achse wird eine komplexe Zahl $f(x)$ zugeordnet. Solche Funktionen
sind den Feldern der Theoretischen Physik~B sehr ähnlich, dem elektrischen Feld oder dem
Potential, nur sind die Werte jetzt komplex. Sie beschreiben die reinen Zustände eines
Teilchens ohne Spin in einer Dimension\zoomref{c03.v2}.

\begin{vertiefung}[id=c03.v2,
    titel={Funktionen als Vektoren},
    wozu={Funktionen lassen sich addieren und mit Zahlen multiplizieren wie Vektoren. Hier
      bekommen sie ein Skalarprodukt, mit einem Integral an der Stelle der Summe (mit einem
      Vorbehalt, der im Fundament aufgelöst wird), und dieselbe Zulassungsbedingung wie die
      Spalten in $\elltwo$.},
    voraussetzung={Abschnitt~\ref{c03.s2}; uneigentliche Integrale. Vertiefung~\ref{c03.v1}
      hilft, ist aber nicht nötig. Der Kern bis Kapitel~4 braucht den Block nicht;
      Kapitel~5 führt den Raum im Kern noch einmal ein.},
    zurueck={Quadratintegrierbare Funktionen bilden einen Vektorraum, auf dem das Integral
      $\braket{f}{g} = \int f^*(x)\, g(x)\, \dd x$ erklärt ist; ein echtes Skalarprodukt wird
      es erst, wenn man Funktionen identifiziert, die fast überall übereinstimmen. Die konstante
      Funktion $1$ gehört nicht dazu.}]
  % Stimme: v03 [31:35]-[32:36].
  Eine Funktion $f\colon \mathbb{R} \to \mathbb{C}$ ordnet jeder Stelle $x$ eine komplexe
  Zahl zu. Das ist grafisch etwas schwierig darzustellen: Für eine komplexe Zahl braucht man
  zwei Achsen, für die Funktion also eine dritte Dimension
  (Abbildung~\ref{c03.fig.komplex}). Funktionen addiert man Stelle für Stelle,
  $(f + g)(x) = f(x) + g(x)$, und ebenso ist $(af)(x) = a f(x)$ für einen Skalar $a$. In
  diesem Sinn ist eine Funktion ein Vektor mit einem Eintrag $f(x)$ für jede Stelle $x$. An
  die Stelle der Summe über $j$ tritt das Integral über $x$.

  \begin{tpcfigure}{c03.fig.komplex}
    \demo{c03.demo.komplex}
    \caption{Eine komplexwertige Funktion braucht drei Dimensionen. Gezeichnet ist
      $f(x) = e^{-x^2/2}\, e^{2ix}$ für $-3 \leq x \leq 3$: Zu jedem $x$ gehört der Punkt
      $(\operatorname{Re} f(x), \operatorname{Im} f(x))$ einer komplexen Ebene, die senkrecht
      auf der $x$-Achse steht. Die Kurve windet sich um die $x$-Achse; ihr Abstand von der
      Achse ist $\abs{f(x)} = e^{-x^2/2}$ (dünne Kreise bei $x = -1, 0, 1$; dünne Speichen verbinden die Achse mit der Kurve). Gestrichelt ist der Realteil
      allein, in der Ebene $\operatorname{Im} f = 0$: die gewohnte Zeichnung, der die
      Hälfte der Information fehlt.}
  \end{tpcfigure}

  In dieser Vertiefung sind alle Funktionen stückweise stetig, so dass jedes Integral ein
  gewöhnliches uneigentliches Integral ist. Was sich ändert, wenn man das fallen lässt,
  steht im Fundament zu diesem Kapitel.

  \tpcdef{D78}

  Das Zeichen $\braket{f}{g}$ steht hier für eine Zahl, so wie $\vb{v}\adj\vb{w}$ in
  $\mathbb{C}^n$. Dieser Kurs schreibt jedes Skalarprodukt so; was der senkrechte Strich
  bedeutet, wird mit der Notation von Dirac klar.

  \begin{proposition}[Quadratintegrierbare Funktionen bilden einen Vektorraum]\label{c03.prop.ltwo}\claim[proof-included]{c03.cl.ltwo}{derived}
    Seien $f$ und $g$ quadratintegrierbar und $a \in \mathbb{C}$. Dann sind $af$ und
    $f + g$ quadratintegrierbar, und das Integral
    $\braket{f}{g} = \int_{-\infty}^{\infty} f^*(x)\, g(x)\, \dd x$ konvergiert absolut, mit
    \begin{equation}\label{c03.eq.ltwo}
      \int_{-\infty}^{\infty} \abs{f^*(x)\, g(x)}\, \dd x
      \leq \tfrac12 \braket{f}{f} + \tfrac12 \braket{g}{g} .
    \end{equation}
    Die Regeln~1 bis~3 von Proposition~\ref{c03.prop.skalarprodukt} gelten sinngemäß, und
    $\braket{f}{f} \geq 0$.
  \end{proposition}
  \begin{proof}
    Für komplexe Zahlen $a, b$ ist $\abs{a^* b} \leq \tfrac12 \abs{a}^2 + \tfrac12 \abs{b}^2$,
    weil $0 \leq (\abs{a} - \abs{b})^2$, und $\abs{a + b}^2 \leq 2\abs{a}^2 + 2\abs{b}^2$, weil
    $\abs{a + b}^2 = \abs{a}^2 + \abs{b}^2 + 2 \operatorname{Re}(a^* b)$ und
    $\operatorname{Re}(a^* b) \leq \abs{a}\,\abs{b}$ ist. Man setzt $a = f(x)$ und $b = g(x)$
    ein; die Ungleichungen gelten dann an jeder Stelle $x$. Integriert über $[-X, X]$ ergeben sie Schranken, die
    nicht von $X$ abhängen. Die Integrale nichtnegativer Funktionen über $[-X, X]$ wachsen
    mit $X$ und sind beschränkt, also konvergieren sie für $X \to \infty$. Ein absolut
    konvergentes uneigentliches Integral konvergiert. Die Regeln~1 bis~3 gelten für den
    Integranden an jeder Stelle, und das Integral ist linear. Schließlich ist
    $\abs{f(x)}^2 \geq 0$ an jeder Stelle.
  \end{proof}

  Regel~4 von Proposition~\ref{c03.prop.skalarprodukt} fehlt mit Absicht in dieser Liste.
  Dass $\braket{f}{f} = 0$ nur für $f = 0$ gilt, ist für Funktionen nicht ganz richtig, und
  die Reparatur steht im Fundament\zoomref{c03.f2}.

  \begin{example}[Ein Nicht-Beispiel]\label{c03.ex.eins}\claim{c03.cl.eins}{derived}
    Die konstante Funktion $f(x) = 1$ ist nicht quadratintegrierbar:
    \begin{equation}\label{c03.eq.eins}
      \int_{-X}^{X} \abs{f(x)}^2\, \dd x = \int_{-X}^{X} 1\, \dd x = 2X
      \;\longrightarrow\; \infty
      \qquad (X \to \infty) .
    \end{equation}
    Die Fläche unter der Kurve ist unendlich groß. Das ist genau das Gegenstück zum Vektor
    mit lauter Einsen aus Beispiel~\ref{c03.ex.einsen}.
  \end{example}

  \begin{example}[Gegen null gehen reicht nicht]\label{c03.ex.gauss}\claim{c03.cl.gauss}{derived}
    Die Funktion $g(x) = e^{-x^2/2}$ ist quadratintegrierbar, und ebenso die Funktion
    $f(x) = e^{-x^2/2}\, e^{2ix}$ aus Abbildung~\ref{c03.fig.komplex}, denn
    $\abs{e^{2ix}} = 1$ und damit $\abs{f(x)}^2 = \abs{g(x)}^2 = e^{-x^2}$. Mit der
    Substitution $x = y/\sqrt{2}$ und dem Gauß-Integral aus Kapitel~\ref{c01},
    $\int_{-\infty}^{\infty} e^{-y^2/2}\, \dd y = \sqrt{2\pi}$, ist
    \begin{equation}\label{c03.eq.gauss}
      \braket{g}{g} = \int_{-\infty}^{\infty} e^{-x^2}\, \dd x
      = \frac{1}{\sqrt{2}} \int_{-\infty}^{\infty} e^{-y^2/2}\, \dd y
      = \sqrt{\pi} .
    \end{equation}
    Die Funktion $h(x) = 1/\sqrt{1 + \abs{x}}$ geht dagegen für $\abs{x} \to \infty$ gegen
    null und ist trotzdem nicht quadratintegrierbar:
    \begin{equation}\label{c03.eq.log}
      \int_{-X}^{X} \frac{\dd x}{1 + \abs{x}} = 2 \int_{0}^{X} \frac{\dd x}{1 + x}
      = 2 \ln(1 + X) \;\longrightarrow\; \infty .
    \end{equation}
    Diese Beispiele zeigen: Dass eine Funktion gegen null geht, reicht nicht. Entscheidend
    ist allein, ob das Integral ihres Betragsquadrats endlich ist.
  \end{example}

  \begin{tpcfigure}{c03.fig.partialintegrale}
    \demo{c03.demo.partialintegrale}
    \caption{Dieselbe Prüfung für Funktionen: das Integral $\int_{-X}^{X} \abs{f(x)}^2\, \dd x$
      als Funktion von $X$. Für $f(x) = e^{-x^2/2}$ (durchgezogen) nähert es sich
      $\sqrt{\pi}$ (dünne Linie, Gleichung~\eqref{c03.eq.gauss}); für
      $f(x) = 1/\sqrt{1 + \abs{x}}$ (gestrichelt) wächst es wie $2\ln(1 + X)$, für $f = 1$
      (strichpunktiert) wie $2X$ (Gleichungen~\eqref{c03.eq.log} und~\eqref{c03.eq.eins}).
      Die Linien zeigen kontinuierlich, was Abbildung~\ref{c03.fig.partialsummen} für den
      diskreten Index $j$ zeigt.}
  \end{tpcfigure}
\end{vertiefung}

Bis einschließlich Kapitel~4 arbeitet der Kern mit endlichdimensionalen Räumen
$\mathbb{C}^n$; für das Qubit ist das $\mathbb{C}^2$. In der Vorlesung stelle ich an dieser
Stelle drei Räume und drei physikalische Systeme nebeneinander, darunter die beiden
unendlichdimensionalen Räume der Vertiefungen\zoomref{c03.v3}.

% Stimme: v03 [37:39]-[39:46]; vl02.tex:138-140.
\begin{vertiefung}[id=c03.v3,
    titel={Drei Räume, drei Systeme},
    wozu={Die drei Räume dieses Kapitels nebeneinander, jeder mit dem physikalischen System,
      zu dem er gehört. Damit sind praktisch alle Hilberträume des Kurses beisammen.},
    voraussetzung={Vertiefungen~\ref{c03.v1} und~\ref{c03.v2}. Der Kern bis Kapitel~4
      braucht den Block nicht.},
    zurueck={Der Kern arbeitet bis einschließlich Kapitel~4 mit $\mathbb{C}^n$; für das
      Qubit ist das $\mathbb{C}^2$.}]
  \begin{greybox}[Drei Räume, drei Systeme]
  Mit $\mathbb{C}^2$ und zwei unendlichdimensionalen Räumen haben wir praktisch alle
  Hilberträume beisammen, die wir in diesem Kurs brauchen werden: $\elltwo$, die Spalten
  mit unendlich vielen Einträgen, deren Betragsquadrate eine endliche Summe haben, und
  $\Ltwo{\mathbb{R}}$, gebildet aus den Funktionen auf der Geraden, deren Betragsquadrat ein
  endliches Integral hat, wobei Funktionen, die fast überall übereinstimmen, als derselbe
  Vektor gelten (Fundament zu diesem Kapitel). Jeder Raum gehört zu einem physikalischen
  System:
  \begin{equation}\label{c03.eq.drei-raeume}\claim{c03.cl.drei-raeume}{source-backed}
    \begin{aligned}
      &\mathbb{C}^2 && \text{\tpcifenglish{qubit}{Qubit}} \\
      &\elltwo && \text{\tpcifenglish{harmonic oscillator}{harmonischer Oszillator}} \\
      &\Ltwo{\mathbb{R}} && \text{\tpcifenglish{spinless particle in one dimension}{Teilchen ohne Spin in einer Dimension}}
    \end{aligned}
  \end{equation}
  Der erste Raum ist endlichdimensional. Die beiden anderen kommen im Kern ab Kapitel~5
  vor.
  \end{greybox}

\begin{tpcfigure}{c03.fig.raeume}
  \demo{c03.demo.raeume}
  \caption{Die drei Räume nebeneinander, jeweils mit den Betragsquadraten eines Vektors.
    (a) $\mathbb{C}^2$: der Vektor $\tfrac12\vb{v}$ mit $\vb{v}$ aus
    Gleichung~\eqref{c03.eq.c2-vektoren}; seine beiden Betragsquadrate sind je $\tfrac12$.
    (b) $\elltwo$: die Betragsquadrate $\abs{v_j}^2 = 1/j^2$ der ersten zwölf Einträge des
    Vektors mit $v_j = 1/j$; es folgen unendlich viele weitere. (c) $\Ltwo{\mathbb{R}}$: das
    Betragsquadrat $\abs{f(x)}^2 = e^{-x^2}$ der Funktion $f(x) = e^{-x^2/2}$. Das
    Skalarprodukt des Vektors mit sich selbst ist links eine Summe von zwei Balken, in der
    Mitte eine Summe von unendlich vielen, rechts die Fläche unter der Kurve.}
  \end{tpcfigure}
\end{vertiefung}

% Stimme: v03 [37:39] („sofern wir noch eine Bedingung ...“), [39:46]-[42:02]; v03 [87:13].
Ein komplexer Vektorraum mit Skalarprodukt ist fast schon ein Hilbertraum. Es fehlt noch eine
Bedingung, die \emph{Vollständigkeit}. Sie ist mathematisch wichtig, aber subtil, und in
diesem Kurs spielt sie keine wesentliche Rolle: In $\mathbb{C}^n$ ist sie immer
erfüllt\zoomref{c03.f3}. Bevor wir den Hilbertraum definieren, brauchen wir noch eine
Schreibweise, die die Rechnungen mit Skalarprodukten sehr viel leichter macht, die Notation
von Dirac. Die eigentliche Hürde am Anfang der Quantenmechanik ist die Notation; da muss man
sich durchkämpfen. Am Anfang ist sie ungewohnt, am Ende macht sie alles leichter, und sie ist
eine Einbahnstraße: Wer sie einmal beherrscht, geht nicht mehr zurück.

% ── Vorlesung 3 (vl03): Dirac-Notation, Skalarprodukte, Hilbertraum, Basen ────────────────
\section{Dirac-Notation}\label{c03.s4}

% Stimme: vl03.tex:5; v03 [43:57]-[44:17] („das ist so, wir müssen damit leben“; „Mathematiker
% mögen diese Notation nicht, Physiker lieben diese Notation“).
Sei $V$ ein komplexer Vektorraum. In der Quantenmechanik wird die folgende Schreibweise
universell verwendet; das ist so, wir müssen damit leben. Mathematiker mögen diese Notation
nicht, Physiker lieben sie: Das ist der Unterschied zwischen Mathematikern und Physikern. Ab
sofort schreiben wir jeden Vektor $\vb{v} \in V$ in der sogenannten Ket-Notation,
\begin{equation*}
  \vb{v} \mapsto \ket{v} .
\end{equation*}
Den Fettdruck lassen wir im Ket weg: Der senkrechte Strich und die spitze Klammer sagen schon,
dass ein Vektor gemeint ist.

% Stimme: vl03.tex:9; v03 [45:35]-[46:10].
Diese Notation bedeutet zunächst nicht mehr als: „$\ket{v}$ steht für einen Vektor im
Vektorraum $V$“. Dass uns die zwei zusätzlichen Symbole einen Vorteil bringen, ist nicht
offensichtlich. Warum wir diese ungewöhnliche Notation eingeführt haben, wird bald klar.

% Stimme: vl03.tex:13; v03 [46:15]-[46:34].
Wir rechnen ein paar Beispiele durch, damit Sie mit der Notation vertraut werden. Es gibt
dabei einige Rechenregeln, die nur implizit in der Notation stecken; man muss sie einmal
gesehen haben.

% Stimme: vl03.tex:15-42; v03 [46:41]-[48:59].
\begin{example}[Kets in $\mathbb{C}^2$]\label{c03.ex.kets}\claim{c03.cl.kets}{derived}
  Sei $V = \mathbb{C}^2$, der Raum, mit dem wir die Zustände eines Qubits beschreiben werden.
  Zwei Beispielvektoren sind
  \begin{equation}\label{c03.eq.e0e1}
    \ket{e_0} = \begin{pmatrix} 1 \\ 0 \end{pmatrix},
    \qquad
    \ket{e_1} = \begin{pmatrix} 0 \\ 1 \end{pmatrix} .
  \end{equation}
  So soll man diese Gleichung lesen: $\ket{e_0}$ steht für diese Spalte, mehr nicht. In einem
  Vektorraum darf man Vektoren addieren und mit Zahlen multiplizieren, und so bilden wir zwei
  neue Vektoren,
  \begin{equation}\label{c03.eq.epm}
    \ket{e_\pm} = \tfrac{1}{\sqrt{2}}\, \ket{e_0} \pm \tfrac{1}{\sqrt{2}}\, \ket{e_1}
    = \frac{1}{\sqrt{2}} \begin{pmatrix} 1 \\ 0 \end{pmatrix} \pm \frac{1}{\sqrt{2}} \begin{pmatrix} 0 \\ 1 \end{pmatrix}
    = \frac{1}{\sqrt{2}} \begin{pmatrix} 1 \\ \pm 1 \end{pmatrix} .
  \end{equation}
  Welcher Vektor dargestellt ist, sieht man, indem man die Spalten aus
  Gleichung~\eqref{c03.eq.e0e1} einsetzt.
\end{example}

% Stimme: v03 [50:02]-[51:04].
% GAP-TJO: In der Vorlesung (v03 [49:20]-[49:46]) folgt der Satz, die Vektoren |e_±> tauchten
% in der Theorie des Drehimpulses, beim Wasserstoffatom und besonders bei Spinsystemen immer
% wieder auf.  Er hat keine schriftliche Quelle (vl03.tex sagt es nicht) und steht deshalb
% nicht im Text (G-0072 (a)).
In den kommenden Kapiteln lassen wir die Spalten meistens weg und rechnen lieber direkt mit
der Notation. Am Anfang sollten Sie aber immer die ursprünglichen Spalten einsetzen, um zu
sehen, was dargestellt wird; in den Übungen werden Sie viele solche Rechnungen durchführen.

% Stimme: vl03.tex:45.  Die zweite Regel (Zahl im Ket) ist eine Ergänzung des Kapitelautors;
% D82 braucht sie.
Es macht übrigens keinen Unterschied, ob wir $\ket{v + w}$ oder $\ket{v} + \ket{w}$
aufschreiben, und ebenso steht $\ket{z v}$ für $z \ket{v}$, wenn $z$ eine Zahl ist:
\begin{equation*}
  \ket{v + w} = \ket{v} + \ket{w} ,
  \qquad
  \ket{z v} = z \ket{v} .
\end{equation*}

% Stimme: vl03.tex:51-66; v03 [52:14]-[53:41] („wenn man diese komischen Symbole sieht, dann
% muss man dieses Objekt als einen Zeilenvektor verstehen“).
Ergänzt wird die Ket-Notation durch die sogenannten \emph{Bra}-Vektoren. Steht $\ket{v}$ für
eine Spalte von $n$ komplexen Zahlen, so macht die Kreuz-Operation aus Abschnitt~\ref{c03.s2}
daraus eine Zeile, und für diese Zeile führen wir ein eigenes Zeichen ein, $\bra{v}$. Wenn Sie
dieses komische Symbol sehen, müssen Sie es als Zeilenvektor verstehen. Bra und Ket zusammen
schreiben das Skalarprodukt.

\tpcdef{D81}

% Stimme: vl03.tex:72-76; v03 [55:09]-[56:11].
\begin{historical}[Bra-Ket]
  Die Notation stammt von Paul Dirac, und sie wird heute universell verwendet
  \srcref{tpc-script-v0.2}{vl03, Kasten zur Schreibweise Bra-Ket}. Englisch \emph{bracket}
  heißt Klammer: Links steht der „bra“, rechts der „ket“, und das „c“ ist irgendwo
  verlorengegangen.
\end{historical}

% Stimme: v03 [56:45]-[59:27] („es gibt keine neue Definition, eine neue Notation für dieselben
% Objekte“; „diese beiden Notationen werden identifiziert“).  Die Einschränkung auf das
% Skalarprodukt aus c03.s2 und der Bra als Zuordnung für andere Skalarprodukte: Reparatur nach
% docs/review/lectures/REVIEW-lec3.md, Zeile SCOPE (de.tex:630), und
% evidence/derivations/c03/REVIEW-lec3.md, B, Zeile SCOPE (de.tex:627).
Das ist keine neue Definition, sondern eine neue Notation für dieselben Objekte und dieselben
Gleichungen. Für dieselbe Zahl haben wir jetzt drei Schreibweisen, die Summe
$\sum_{j} v_j^* w_j$, das Produkt $\vb{v}\adj\vb{w}$ und das Bra-Ket $\braket{v}{w}$, und
ausgerechnet wird ein Bra-Ket nach der alten Regel. Zwischen $\braket{v}{w}$ und
$(\ket{v})\adj \ket{w}$ machen wir gar keinen Unterschied mehr. Das gilt für das Skalarprodukt
aus Abschnitt~\ref{c03.s2}, und mit ihm rechnen wir weiterhin, wenn wir Spalten einsetzen. Im
nächsten Abschnitt kommen auch andere Skalarprodukte vor. Für sie ist der Bra $\bra{v}$ die
Zuordnung, die jedem Ket $\ket{w}$ die Zahl $\braket{v}{w}$ zuordnet, und im Allgemeinen nicht
mehr die konjugierte Zeile (Beispiel~\ref{c03.ex.gewichtet}).

% Stimme: vl03.tex:78-87; v03 [59:37]-[61:04].
\begin{example}[Ein Bra-Ket in $\mathbb{C}^2$]\label{c03.ex.braket}\claim{c03.cl.braket}{derived}
  Mit den Vektoren aus Beispiel~\ref{c03.ex.kets} ist
  \begin{equation}\label{c03.eq.e0eplus}
    \braket{e_0}{e_+} = (\ket{e_0})\adj \ket{e_+}
    = \begin{pmatrix} 1 & 0 \end{pmatrix}
      \begin{pmatrix} \frac{1}{\sqrt{2}} \\ \frac{1}{\sqrt{2}} \end{pmatrix}
    = \frac{1}{\sqrt{2}} .
  \end{equation}
\end{example}

% Stimme: vl03.tex:89-93; v03 [62:30]-[65:06].
Damit die Notation ihren Vorteil ausspielen kann, braucht man eine Regel dafür, wie sich die
Kreuz-Operation auf Linearkombinationen auswirkt. Für komplexe Zahlen $\alpha, \beta$ und
Kets $\ket{v}, \ket{w}$ gilt
\begin{equation}\label{c03.eq.bra-regel}\claim{c03.cl.bra-regel}{derived}
  (\alpha \ket{v} + \beta \ket{w})\adj = \alpha^* \bra{v} + \beta^* \bra{w} ,
\end{equation}
denn die Kreuz-Operation konjugiert jeden Eintrag, und
$(\alpha v_j + \beta w_j)^* = \alpha^* v_j^* + \beta^* w_j^*$. Genau diese Regel macht die
Notation nützlich: Mit ihr rechnen wir Skalarprodukte aus, ohne die Vektoren je als Spalten
hinzuschreiben. In $\mathbb{C}^2$ ist das Hinschreiben nicht schlimm. In unendlichdimensionalen
Räumen kann man einen Vektor aber im Allgemeinen gar nicht vollständig hinschreiben, und dann
ist die Rechnung mit Kets und dieser Regel der Weg, effizient zu rechnen.

% Stimme: vl03.tex:95-125; v03 [65:32]-[71:18].
\begin{example}[Ein Skalarprodukt nur mit Kets]\label{c03.ex.dirac}\claim{c03.cl.dirac}{derived}
  Gesucht ist das Skalarprodukt $\braket{v}{w}$ von
  \begin{equation}\label{c03.eq.dirac-vw}
    \ket{v} = \begin{pmatrix} 1 \\ i \end{pmatrix},
    \qquad
    \ket{w} = \begin{pmatrix} 1 \\ 1 \end{pmatrix},
  \end{equation}
  diesmal ohne die Spalten einzusetzen. Zuerst drücken wir beide Vektoren durch die Kets aus
  Beispiel~\ref{c03.ex.kets} aus: $\ket{v} = \ket{e_0} + i \ket{e_1}$ und
  $\ket{w} = \ket{e_0} + \ket{e_1}$. Nach Gleichung~\eqref{c03.eq.bra-regel} ist
  $\bra{v} = \bra{e_0} - i \bra{e_1}$; beachten Sie das Minuszeichen. Ausmultiplizieren gibt
  vier Skalarprodukte:
  \begin{equation}\label{c03.eq.dirac}
    \begin{aligned}
      \braket{v}{w} &= (\bra{e_0} - i \bra{e_1}) (\ket{e_0} + \ket{e_1}) \\
      &= \braket{e_0}{e_0} + \braket{e_0}{e_1} - i \braket{e_1}{e_0} - i \braket{e_1}{e_1} \\
      &= 1 + 0 - 0 - i = 1 - i ,
    \end{aligned}
  \end{equation}
  denn $\braket{e_0}{e_0} = \braket{e_1}{e_1} = 1$ und $\braket{e_0}{e_1} = \braket{e_1}{e_0} = 0$.
  Dasselbe Ergebnis hätten wir bekommen, wenn wir einfach die Spalten eingesetzt hätten:
  $\braket{v}{w} = \begin{pmatrix} 1 & -i \end{pmatrix} \begin{pmatrix} 1 \\ 1 \end{pmatrix} = 1 - i$.
\end{example}

% Stimme: vl03.tex:127; v03 [71:22]-[72:54].  Die Dimensionszahl aus Skript und Vorlesung
% („ab fünf oder sechs Dimensionen“, v03: „ab 5, 6, 7 Dimensionen“) ist weggelassen: Sie ist
% eine Einschätzung ohne Herleitung (COORDINATION.md Regel 10); siehe vl03-map.json, text@127.
Seien wir ehrlich: In zwei Dimensionen ist der erste Weg aufwendiger. In höheren Dimensionen
werden Sie aber ganz schnell eine Liebe für die Ket-Notation entdecken. Ich verwende sie
trotzdem auch in zwei Dimensionen, weil sie sehr schön mit Operatoren und Matrizen
zusammenspielt; das sehen wir bei den Operatoren.

% Stimme: v03 [73:22]-[73:44].
Sobald Sie die Ket-Notation beherrschen, können Sie die Gleichungen der Quantenmechanik lesen.
Die Übersetzung ist immer dieselbe: Ket steht für Vektor.

\section{Skalarprodukte}\label{c03.s5}

% Stimme: vl03.tex:130-132; v03 [73:51]-[75:47] („für immer und ewig“; „Skalarprodukt ist nicht
% kanonisch“).  Der zweite und dritte Satz nach docs/review/lectures/REVIEW-lec3.md, Zeile
% MISMATCH (de.tex:700): Die Vertiefungen c03.v1 und c03.v2 haben schon andere Skalarprodukte.
Bevor wir den Hilbertraum definieren können, müssen wir das Konzept des Skalarprodukts
sauberer und mathematisch besser fassen. In $\mathbb{C}^n$ haben wir bisher immer das
Skalarprodukt $\vb{v}\adj\vb{w}$ verwendet. Aber selbst auf diesem Vektorraum gibt es andere
Skalarprodukte, sogar unendlich viele (Beispiel~\ref{c03.ex.gewichtet}). In der Physik wählt man normalerweise eines
aus und arbeitet dann mit diesem Skalarprodukt für immer und ewig; mathematisch gibt es
mehrere Möglichkeiten, und manchmal ist es hilfreich, ein bisschen flexibel zu sein. Ein
komplexer Vektorraum kann auch ganz ohne Skalarprodukt existieren: Das Skalarprodukt ist eine
Zusatzwahl, keine Struktur, die der Vektorraum von selbst mitbringt. Wir brauchen eines, weil
die Quantenmechanik probabilistisch ist; es ist unser Werkzeug für Wahrscheinlichkeiten.

% Stimme: vl03.tex:134-147; v03 [76:10]-[82:34] („nicht jede solche Zuordnung zählt als
% Skalarprodukt“).
Ein Skalarprodukt ordnet je zwei Vektoren eine komplexe Zahl zu, die wir in der Notation von
Dirac schreiben. Nicht jede solche Zuordnung zählt als Skalarprodukt; sie muss ein paar
Bedingungen erfüllen. Diese Bedingungen beschreiben auch die anderen Skalarprodukte. Für die
Rechnungen mit Spalten verwenden wir weiterhin das Skalarprodukt aus Abschnitt~\ref{c03.s2}.

\tpcdef{D82}

% Stimme: vl03.tex:139; v03 [79:26]-[80:06].
Im ersten Argument kommt beim Herausziehen einer Zahl also ein Stern dazu. Diese Regeln
gelten im endlichdimensionalen wie im unendlichdimensionalen Fall; deshalb sind Notation und
Begriff so wichtig. Das Skalarprodukt $\vb{v}\adj\vb{w}$ von $\mathbb{C}^n$ erfüllt alle drei
Bedingungen: Das sind die Regeln~2 bis~4 von Proposition~\ref{c03.prop.skalarprodukt}. Die
Komplexkonjugation in der zweiten Regel für Zahlen ist wichtig, damit $\braket{\psi}{\psi}$
positiv sein kann. Wäre das Skalarprodukt auch im ersten Argument linear, so wäre für jeden
Vektor
\begin{equation}\label{c03.eq.warum-stern}\claim{c03.cl.warum-stern}{derived}
  \braket{i\psi}{i\psi} = i \cdot i\, \braket{\psi}{\psi} = -\braket{\psi}{\psi} ,
\end{equation}
und beide Seiten könnten nur dann $\geq 0$ sein, wenn $\braket{\psi}{\psi} = 0$ wäre, auch für
$\ket{\psi} \neq 0$. Beispiel~\ref{c03.ex.ohne-stern} zeigt dasselbe an einem einzelnen
Vektor.

\begin{example}[Unendlich viele Skalarprodukte auf $\mathbb{C}^2$]\label{c03.ex.gewichtet}\claim{c03.cl.gewichtet}{derived}
  Sei $c > 0$ eine reelle Zahl. Für $\vb{v}, \vb{w} \in \mathbb{C}^2$ setzen wir
  \begin{equation}\label{c03.eq.gewichtet}
    \braket{v}{w}_c = v_1^* w_1 + c\, v_2^* w_2 .
  \end{equation}
  Das ist für jedes $c > 0$ ein Skalarprodukt: Die Additivität und die beiden Regeln für Zahlen
  gelten Summand für Summand, genau wie in Proposition~\ref{c03.prop.skalarprodukt}, und
  $\braket{v}{v}_c = \abs{v_1}^2 + c \abs{v_2}^2$ ist $\geq 0$ und nur dann null, wenn
  $v_1 = v_2 = 0$ ist. Für $c = 1$ ist es das Skalarprodukt von $\mathbb{C}^2$. Für verschiedene
  $c$ sind es verschiedene Skalarprodukte, denn $\braket{e_1}{e_1}_c = c$. Es gibt also
  unendlich viele. Zu jedem gehört ein eigener Bra: Wegen
  $\braket{v}{w}_c = \begin{pmatrix} v_1^* & c\, v_2^* \end{pmatrix} \vb{w}$ steht $\bra{v}$
  hier für die Zeile $\begin{pmatrix} v_1^* & c\, v_2^* \end{pmatrix}$, und das ist für
  $c \neq 1$ und $v_2 \neq 0$ nicht die Zeile $\vb{v}\adj$.
\end{example}

% Stimme: vl03.tex:148-151; v03 [82:34]-[83:43] („wird oft auch einfach Betrag von psi
% geschrieben“).
Die Positivität erlaubt uns, die Länge eines Vektors zu definieren.

\tpcdef{D83}

In $\mathbb{C}^n$ mit dem Skalarprodukt aus Abschnitt~\ref{c03.s2} ist
$\norm{v}^2 = \vb{v}\adj\vb{v} = \sum_j \abs{v_j}^2$ (Gleichung~\eqref{c03.eq.selbst}). Statt $\norm{\psi}$ schreibt man oft auch einfach „Betrag
von $\psi$“.

% Stimme: vl03.tex:152-155; v04 [02:20]-[03:31] („diese Positivitätsbedingungen sind sehr
% stark, die haben sehr schöne Folgen“; „das ist nichts Neues“).
Die Positivität ist eine sehr starke Bedingung, und sie hat zwei sehr schöne Folgen. Die erste
ist die konjugierte Symmetrie. Für die Skalarprodukte, die wir bisher gesehen haben, ist sie
nichts Neues; hier folgt sie aus den drei Bedingungen allein.

\begin{proposition}[Konjugierte Symmetrie]\label{c03.prop.symmetrie}\claim[proof-included]{c03.cl.symmetrie}{derived}
  Für jedes Skalarprodukt und alle Vektoren $\ket{\phi}, \ket{\psi}$ gilt
  \begin{equation}\label{c03.eq.symmetrie}
    \braket{\phi}{\psi} = \braket{\psi}{\phi}^* .
  \end{equation}
\end{proposition}
\begin{proof}
  Nach der Positivität ist $\braket{\xi}{\xi}$ für jeden Vektor $\ket{\xi}$ eine reelle
  Zahl. Sei $a = \braket{\phi}{\psi}$ und $b = \braket{\psi}{\phi}$. Mit der Additivität ist
  \begin{equation}\label{c03.eq.symmetrie-beweis}
    \braket{\phi + \psi}{\phi + \psi} = \braket{\phi}{\phi} + \braket{\psi}{\psi} + a + b ,
  \end{equation}
  also ist $a + b$ reell. Dieselbe Rechnung mit $i\ket{\psi}$ statt $\ket{\psi}$: Nach den
  Regeln für Zahlen ist $\braket{\phi}{i\psi} = i a$, $\braket{i\psi}{\phi} = -i b$ und
  $\braket{i\psi}{i\psi} = (-i)\, i\, \braket{\psi}{\psi} = \braket{\psi}{\psi}$, also ist
  $i a - i b = i (a - b)$ reell. Aus $\operatorname{Im}(a + b) = 0$ folgt
  $\operatorname{Im} b = -\operatorname{Im} a$, und weil $a - b$ rein imaginär ist, ist
  $\operatorname{Re} b = \operatorname{Re} a$. Also ist $b = a^*$.
\end{proof}

Reed und Simon nehmen die konjugierte Symmetrie als vierte Bedingung in die Definition des
Skalarprodukts auf \srcref{reed-simon-1}{S.~36}; nach der Proposition ist das dieselbe
Definition.

% Stimme: vl03.tex:156-159; v04 [03:38]-[05:23] („besonders nützlich in der Quantenmechanik“;
% „gilt für alle phi und psi im Hilbertraum“).
Die zweite Folge ist neu, die \emph{Schwarz'sche Ungleichung}. Sie ist besonders schön und in
der Quantenmechanik besonders nützlich.

\begin{proposition}[Schwarz'sche Ungleichung]\label{c03.prop.schwarz}\claim[proof-included]{c03.cl.schwarz}{derived}
  Für jedes Skalarprodukt und alle Vektoren $\ket{\phi}, \ket{\psi}$ gilt
  \begin{equation}\label{c03.eq.schwarz}
    \abs{\braket{\phi}{\psi}} \leq \norm{\phi}\, \norm{\psi} .
  \end{equation}
\end{proposition}
\begin{proof}
  Für $\ket{\psi} = 0$ sind beide Seiten null. Sonst ist $\norm{\psi}^2 > 0$. Wir setzen
  $c = \braket{\psi}{\phi} / \norm{\psi}^2$ und $\ket{\xi} = \ket{\phi} - c \ket{\psi}$. Nach
  Proposition~\ref{c03.prop.symmetrie} ist $c^* = \braket{\phi}{\psi} / \norm{\psi}^2$. Mit der
  Additivität, den Regeln für Zahlen und der Positivität für $\ket{\xi}$ ist
  \begin{equation}\label{c03.eq.schwarz-beweis}
    0 \leq \braket{\xi}{\xi}
      = \norm{\phi}^2 - c\, \braket{\phi}{\psi} - c^*\, \braket{\psi}{\phi} + c^* c\, \norm{\psi}^2
      = \norm{\phi}^2 - \frac{\abs{\braket{\phi}{\psi}}^2}{\norm{\psi}^2} ,
  \end{equation}
  denn bis auf das Vorzeichen ist jeder der drei letzten Summanden gleich dem Bruch auf der
  rechten Seite. Multiplikation mit $\norm{\psi}^2$ und die Wurzel geben die Behauptung.
\end{proof}

% Stimme: vl03.tex:160-163; v04 [05:32]-[06:20].
Mit der Schwarz'schen Ungleichung kann man eine weitere, ebenfalls sehr nützliche Ungleichung
beweisen, die Dreiecksungleichung: Die Norm einer Summe von Vektoren ist nie größer als die
Summe der Normen.

\begin{proposition}[Dreiecksungleichung]\label{c03.prop.dreieck}\claim[proof-included]{c03.cl.dreieck}{derived}
  Für jedes Skalarprodukt und endlich viele Vektoren $\ket{\psi_1}, \dots, \ket{\psi_N}$ gilt
  \begin{equation}\label{c03.eq.dreieck}
    \norm{\sum_{j=1}^{N} \ket{\psi_j}} \leq \sum_{j=1}^{N} \norm{\psi_j} .
  \end{equation}
\end{proposition}
\begin{proof}
  Für zwei Vektoren ist nach der Additivität und Proposition~\ref{c03.prop.symmetrie}
  \begin{equation}\label{c03.eq.dreieck-beweis}
    \norm{\phi + \psi}^2
      = \norm{\phi}^2 + \norm{\psi}^2 + 2 \operatorname{Re} \braket{\phi}{\psi}
      \leq \norm{\phi}^2 + \norm{\psi}^2 + 2 \norm{\phi}\, \norm{\psi}
      = \bigl( \norm{\phi} + \norm{\psi} \bigr)^2 ,
  \end{equation}
  denn $\braket{\psi}{\phi} = \braket{\phi}{\psi}^*$, $\operatorname{Re} z \leq \abs{z}$ für
  jede komplexe Zahl $z$, und dann die Schwarz'sche Ungleichung. Die Wurzel gibt
  $\norm{\phi + \psi} \leq \norm{\phi} + \norm{\psi}$. Für $N$ Vektoren wendet man das
  wiederholt an: $\norm{\psi_1 + \dots + \psi_N} \leq \norm{\psi_1 + \dots + \psi_{N-1}} + \norm{\psi_N}$,
  und so weiter bis zu einem einzigen Vektor.
\end{proof}

Reed und Simon zeigen auf demselben Weg, dass die Norm eines Vektorraums mit Skalarprodukt
eine Norm im Sinne der normierten Räume ist \srcref{reed-simon-1}{S.~38}.

% Stimme: vl03.tex:164; v04 [04:28]-[04:45], [06:37]-[06:45], [09:26]-[10:11] („diese
% Eigenschaften werden wir quasi ständig verwenden“; „gekoppelt mit der Heisenbergschen
% Ungleichung“).  Nicht übernommen: das Beispiel aus v04 („um die Energie eines Systems zu
% beschränken“); weder vl03.tex:164 noch eine Herleitung gibt eine solche Schranke
% (evidence/derivations/c03/REVIEW-lec3.md, B, Zeile UNSUPPORTED).
Beide Ungleichungen sind Arbeitsgeräte, die wir quasi ständig verwenden. Die wichtigste ist
die Schwarz'sche Ungleichung: Sie steckt hinter der Heisenbergschen Unschärferelation
\srcref{tpc-script-v0.2}{vl03, Abschnitt Skalarprodukte}, die in Kapitel~6 kommt. Mit der
Dreiecksungleichung schätzen wir Größen ab.

% Stimme: vl03.tex:166-170; v04 [06:51]-[08:44].
Eine letzte Definition, bevor wir die Skalarprodukte verlassen.

\tpcdef{D84}

Das ist genau analog zu den gewöhnlichen Vektoren im dreidimensionalen Raum. Achtung: Ein
verschwindendes Skalarprodukt heißt \emph{nicht}, dass einer der Vektoren null ist. Das
erzwingt die Positivität nur im Fall $\ket{\phi} = \ket{\psi}$.

\begin{example}[Orthonormale Vektoren in $\mathbb{C}^2$]\label{c03.ex.orthogonal}\claim{c03.cl.orthogonal}{derived}
  Die Vektoren $\ket{e_\pm}$ aus Beispiel~\ref{c03.ex.kets} sind nicht null, und doch ist
  \begin{equation}\label{c03.eq.orthogonal}
    \braket{e_+}{e_-} = \frac{1}{2} \begin{pmatrix} 1 & 1 \end{pmatrix} \begin{pmatrix} 1 \\ -1 \end{pmatrix} = 0 ,
    \qquad
    \norm{e_\pm}^2 = \frac{1}{2} \bigl( 1 + (\pm 1)^2 \bigr) = 1 .
  \end{equation}
  Ebenso ist $\braket{e_0}{e_1} = 0$ und $\norm{e_0} = \norm{e_1} = 1$. Beide Paare sind
  orthonormal.
\end{example}

% Stimme: v04 [08:47]-[09:22].
Man darf die Norm auch als Länge eines Vektors lesen. Die Vektoren der Quantenmechanik leben
in einem abstrakten Raum, der mit unseren drei Dimensionen gar nichts zu tun hat; trotzdem ist
es manchmal hilfreich, ihre Normen als Längen und orthogonale Vektoren als senkrecht
aufeinander zu lesen.

\section{Hilbertraum}\label{c03.s6}

% Stimme: v03 [83:59]-[84:03], [86:43]-[87:09]; vl03.tex:176-178.
Jetzt haben wir genug beisammen, um Hilberträume zu definieren.

\tpcdef{D14}

% Stimme: vl03.tex:177 (Fußnote), 180; v03 [87:13]-[88:24], [89:12]-[89:21] („ein sehr
% wichtiger, aber subtile mathematische Bedingung“; „diese Bedingung hat mit Konvergenz in
% unendlichdimensionalen Räumen zu tun“; „ein Hilbertraum ist nichts anderes als ein komplexer
% Vektorraum mit Skalarprodukt“).  Was „vollständig“ heißt, steht seit der Reparatur in D14; die
% Abkürzung gilt nur für die Rechnungen in C^n mit dem Skalarprodukt aus c03.s2
% (docs/review/lectures/REVIEW-lec3.md, Zeile SCOPE (de.tex:878);
% evidence/derivations/c03/REVIEW-lec3.md, B, Zeile SCOPE (de.tex:883)).  Dass die Entwicklung
% nach einer unendlichen Basis die Vollständigkeit benutzt: Satz c03.thm.parseval setzt einen
% Hilbertraum voraus, und der Beweis bei Reed-Simon I S. 45 benutzt sie für den Grenzwert.
Die Vollständigkeit ist eine sehr wichtige, aber subtile mathematische Bedingung. Sie hat mit
Konvergenz in unendlichdimensionalen Räumen zu tun, und in unserer Vorlesung spielt sie keine
wesentliche Rolle. Für unsere Rechnungen in $\mathbb{C}^n$ mit dem Skalarprodukt aus
Abschnitt~\ref{c03.s2} ist sie automatisch erfüllt\zoomref{c03.f3}. Dort können wir alle
Rechnungen vereinfachen, indem wir sie ignorieren, und für diese Rechnungen ist ein
Hilbertraum schlicht ein komplexer Vektorraum mit Skalarprodukt. Der Raum $\mathbb{C}^n$ ist
also ein Hilbertraum, und für das Qubit ist das $\mathbb{C}^2$. In unendlich vielen
Dimensionen gilt die Abkürzung nicht: Dort benutzt man die Vollständigkeit, etwa für die
Entwicklung nach einer Basis mit unendlich vielen Vektoren.

\section{Basen}\label{c03.s7}

% Stimme: vl03.tex:184; v04 [10:27]-[11:24].
Das nächste Konzept in der Theorie der Hilberträume ist die \emph{Basis}. Eine Basis ist eine
Menge von Vektoren mit sehr schönen Eigenschaften: Nach ihr kann man alle anderen Vektoren des
Hilbertraums zerlegen. Wir brauchen Basen aus einem praktischen Grund: Die meisten Rechnungen
der Quantenmechanik werden erst dann explizit, wenn eine Basis gewählt ist. Normalerweise
wählt man eine Basis aus und macht alle Rechnungen bezüglich dieser Basis.

% Stimme: vl03.tex:186-190; v04 [11:32]-[16:26].  Endliche Indexmenge: Entscheidung D4.
Im Kern ist die Indexmenge $X$ einer Basis endlich.

\tpcdef{D85}

% Stimme: vl03.tex:191; v04 [14:22]-[15:16] („und verwirrenderweise sagen wir, dass diese
% Basis vollständig ist“).  „ein anderer Begriff“ statt „gar nichts zu tun“: Reed-Simon I S. 45
% benutzt die Vollständigkeit des Raums im Beweis der Entwicklung nach einer Basis
% (evidence/derivations/c03/REVIEW-lec3.md, B, Zeile MINOR (de.tex:907)).
Zwei Warnungen. Verwirrenderweise heißt die zweite Eigenschaft \emph{Vollständigkeit}; das
ist ein anderer Begriff von Vollständigkeit als die des Hilbertraums aus
Abschnitt~\ref{c03.s6}. Und
$\kron{j}{k}$ ist das Kronecker-Delta, nicht die Diracsche Delta-Funktion, die später
vorkommt.

% Stimme: v04 [16:55]-[22:29].
\begin{example}[Zwei Basen von $\mathbb{C}^2$]\label{c03.ex.basis}\claim{c03.cl.basis}{derived}
  Die Vektoren $\ket{e_0}, \ket{e_1}$ bilden eine Basis von $\mathbb{C}^2$, mit der Indexmenge
  $X = \{0, 1\}$. Orthonormal sind sie nach Beispiel~\ref{c03.ex.orthogonal}; zu prüfen waren
  hier vier Skalarprodukte. Die Vollständigkeit ist etwas anspruchsvoller. Jeder Vektor von
  $\mathbb{C}^2$ lässt sich als $\ket{\psi} = \alpha \ket{e_0} + \beta \ket{e_1}$ schreiben, mit
  den Einträgen $\alpha, \beta \in \mathbb{C}$, und wir rechnen beide Seiten der
  Vollständigkeitsbedingung getrennt aus. Links steht
  \begin{equation}\label{c03.eq.basis-links}
    \norm{\psi}^2 = (\alpha^* \bra{e_0} + \beta^* \bra{e_1}) (\alpha \ket{e_0} + \beta \ket{e_1})
    = \abs{\alpha}^2 + \abs{\beta}^2 ,
  \end{equation}
  rechts, mit $\braket{e_0}{\psi} = \alpha$ und $\braket{e_1}{\psi} = \beta$,
  \begin{equation}\label{c03.eq.basis-rechts}
    \abs{\braket{e_0}{\psi}}^2 + \abs{\braket{e_1}{\psi}}^2 = \abs{\alpha}^2 + \abs{\beta}^2 .
  \end{equation}
  Beide Seiten stimmen überein. Genauso bilden $\ket{e_+}, \ket{e_-}$ eine Basis: Sie sind
  orthonormal (Beispiel~\ref{c03.ex.orthogonal}), und mit
  $\braket{e_\pm}{\psi} = (\alpha \pm \beta) / \sqrt{2}$ ist
  \begin{equation}\label{c03.eq.basis-pm}
    \abs{\braket{e_+}{\psi}}^2 + \abs{\braket{e_-}{\psi}}^2
    = \tfrac12 \bigl( \abs{\alpha + \beta}^2 + \abs{\alpha - \beta}^2 \bigr)
    = \abs{\alpha}^2 + \abs{\beta}^2 .
  \end{equation}
\end{example}

% Stimme: v04 [18:07]-[18:27], [22:59]-[23:39] („wenn Sie ein bisschen Langeweile haben“;
% „normalerweise ... bildet dieses System gar keine Basis“).
Bei $n$ Basisvektoren sind für die Orthonormalität $n^2$ Gleichungen zu prüfen; wegen der
konjugierten Symmetrie genügen die mit $j \leq k$, und das sind immer noch eine ganze Menge.
Wenn Sie ein bisschen Langeweile haben, schreiben Sie ein paar Vektoren auf und prüfen Sie, ob
sie eine Basis bilden. Wählt man die Vektoren beliebig aus, so bilden sie im Allgemeinen
keine Basis: Schon $\ket{e_0}$ und $\ket{e_+}$ sind nach Beispiel~\ref{c03.ex.braket} nicht
orthogonal.

% Stimme: vl03.tex:193-201; v04 [24:37]-[25:48] („sieht harmlos aus, aber wir wenden diese
% Gleichung ständig an“).
Die wichtigste Eigenschaft einer Basis ist, dass man jeden Vektor nach ihr entwickeln kann.

\begin{proposition}[Entwicklung nach einer Basis]\label{c03.prop.entwicklung}\claim[proof-included]{c03.cl.entwicklung}{derived}
  Sei $\{\ket{\phi_j}\}_{j \in X}$ eine Basis von $\Hil$ mit endlicher Indexmenge $X$. Dann
  gilt für jeden Vektor $\ket{\Psi} \in \Hil$
  \begin{equation}\label{c03.eq.vollstaendigkeit}
    \ket{\Psi} = \sum_{j \in X} \braket{\phi_j}{\Psi}\, \ket{\phi_j} .
  \end{equation}
  Umgekehrt ist ein orthonormales System mit endlicher Indexmenge, nach dem sich jeder Vektor
  so entwickeln lässt, vollständig, also eine Basis.
\end{proposition}
\begin{proof}
  Sei $c_j = \braket{\phi_j}{\Psi}$ und $\ket{\xi} = \ket{\Psi} - \sum_j c_j \ket{\phi_j}$.
  Wegen der Orthonormalität ist $\braket{\phi_k}{\xi} = c_k - \sum_j c_j \kron{k}{j} = 0$ für
  jedes $k$. Mit den Regeln für Zahlen im ersten Argument folgt
  $\braket{\xi}{\xi} = \braket{\Psi}{\xi} - \sum_j c_j^* \braket{\phi_j}{\xi} = \braket{\Psi}{\xi}$,
  und nach Proposition~\ref{c03.prop.symmetrie} ist $\braket{\Psi}{\phi_j} = c_j^*$, also
  \begin{equation}\label{c03.eq.entwicklung-beweis}
    \braket{\xi}{\xi} = \braket{\Psi}{\Psi} - \sum_{j} c_j\, \braket{\Psi}{\phi_j}
    = \norm{\Psi}^2 - \sum_{j} \abs{c_j}^2 = 0
  \end{equation}
  nach der Vollständigkeit. Nach der Positivität ist $\ket{\xi} = 0$, und das ist
  Gleichung~\eqref{c03.eq.vollstaendigkeit}. Umgekehrt: Gilt die Entwicklung, so ist
  $\norm{\Psi}^2 = \sum_{j,k} c_j^* c_k \braket{\phi_j}{\phi_k} = \sum_j \abs{c_j}^2$.
\end{proof}

\begin{merksatz}[Vollständigkeitsrelation]
  Die Gleichung
  $\ket{\Psi} = \sum_{j \in X} \braket{\phi_j}{\Psi}\, \ket{\phi_j}$ ist eine Form der
  sogenannten \emph{Vollständigkeitsrelation}: Jeder Vektor ist die Summe seiner Komponenten
  längs der Basisvektoren, und die Koeffizienten sind Skalarprodukte. Sie sieht harmlos aus,
  wird aber in der Quantenmechanik ständig angewandt.
\end{merksatz}

% Stimme: vl03.tex:205-210 (Abbildungstext des Skripts).
\begin{example}[Ein Vektor, zwei Basen]\label{c03.ex.zwei-basen}\claim{c03.cl.zwei-basen}{derived}
  Sei $\ket{\Psi} = \tfrac{1}{\sqrt{5}} \bigl( 2 \ket{e_0} + \ket{e_1} \bigr)$. In der Basis
  $\ket{e_0}, \ket{e_1}$ sind die Koeffizienten $c_0 = \braket{e_0}{\Psi} = 2/\sqrt{5}$ und
  $c_1 = \braket{e_1}{\Psi} = 1/\sqrt{5}$, in der Basis $\ket{e_+}, \ket{e_-}$ sind es
  $c_\pm = \braket{e_\pm}{\Psi} = (2 \pm 1)/\sqrt{10}$. Die Gewichte
  \begin{equation}\label{c03.eq.gewichte}
    \abs{c_0}^2 = \tfrac45 , \quad \abs{c_1}^2 = \tfrac15 ,
    \qquad
    \abs{c_+}^2 = \tfrac{9}{10} , \quad \abs{c_-}^2 = \tfrac{1}{10}
  \end{equation}
  sind in den beiden Basen verschieden; ihre Summe ist in beiden Fällen $1 = \norm{\Psi}^2$.
\end{example}

Abbildung~\ref{c03.fig.vollst} rechnet das an den beiden Basen von $\mathbb{C}^2$ vor. In (a)
schließt sich jede der beiden Zerlegungen desselben $\ket{\Psi}$ zu einem eigenen Rechteck:
verschiedene Seiten, dieselbe Diagonale. In (b) steht, was von der Basis unabhängig bleibt:
nicht die einzelnen Gewichte, sondern ihre Summe.

\begin{tpcfigure}{c03.fig.vollst}
  \demo{c03.demo.vollst}
  \caption{Die Vollständigkeitsrelation an einem Vektor in zwei Basen. (a) Der Vektor
    $\ket{\Psi} = (2\ket{e_0} + \ket{e_1})/\sqrt{5}$ aus Beispiel~\ref{c03.ex.zwei-basen}
    (dicker Pfeil) als Summe seiner Komponenten $c_j \ket{\phi_j}$ mit
    $c_j = \braket{\phi_j}{\Psi}$, einmal in der Basis $\ket{e_0}, \ket{e_1}$ (Rechteck
    gestrichelt geschlossen), einmal in der Basis $\ket{e_+}, \ket{e_-}$ (Rechteck gepunktet
    geschlossen). Weil alle Koeffizienten reell sind, liegt alles in der reellen Ebene, die
    $\ket{e_0}$ und $\ket{e_1}$ aufspannen; die Strahlen von $\ket{e_\pm}$ halbieren die Winkel
    zwischen den Achsen. Die beiden Rechtecke haben verschiedene Seiten und dieselbe
    Diagonale. (b) Die Gewichte $\abs{c_j}^2$ als Balkenstücke: $\tfrac45$ und $\tfrac15$
    gegen $\tfrac{9}{10}$ und $\tfrac{1}{10}$. Die Aufteilung hängt von der Basis ab, die Summe
    $\norm{\Psi}^2 = 1$ (gestrichelte Linie) nicht.}
\end{tpcfigure}

% Stimme: vl03.tex:212; v04 [25:56]-[27:16] („es gibt immer eine Basis in endlichen
% Dimensionen“; „die Existenz von einer Basis ist noch nicht klar, man muss das beweisen, man
% weiß, dass es sehr schwer zu beweisen ist“).
In endlich vielen Dimensionen ist es relativ leicht, eine Basis zu finden. In $\mathbb{C}^n$
mit dem Skalarprodukt aus Abschnitt~\ref{c03.s2} nimmt man einfach die Spalten $\ket{u_1}, \dots, \ket{u_n}$, wobei $\ket{u_j}$ im $j$-ten
Eintrag eine $1$ hat und sonst lauter Nullen; für $n = 2$ sind das $\ket{e_0}$ und
$\ket{e_1}$. Für einen Vektor $\ket{\psi}$ mit den Einträgen $\psi_1, \dots, \psi_n$ ist
\begin{equation}\label{c03.eq.kanonisch}\claim{c03.cl.kanonisch}{derived}
  \braket{u_j}{u_k} = \kron{j}{k} ,
  \qquad
  \sum_{j=1}^{n} \abs{\braket{u_j}{\psi}}^2 = \sum_{j=1}^{n} \abs{\psi_j}^2 = \norm{\psi}^2 .
\end{equation}
Im unendlichdimensionalen Fall ist alles schwieriger: Dass es eine Basis gibt, ist nicht klar,
man muss es beweisen, und der Beweis ist schwer\zoomref{c03.v4}.

\begin{vertiefung}[id=c03.v4,
    titel={Basen mit unendlich vielen Vektoren},
    wozu={Die Basen von $\elltwo$ und $\Ltwo{\mathbb{R}}$ in diesem Block haben unendlich
      viele Vektoren, und die Summen der Vollständigkeitsrelation werden Reihen. Hier steht,
      wie man eine Basis dann definiert, dass es immer eine gibt und dass die Entwicklung nach
      einer Basis mit abzählbarer Indexmenge gilt, mit einer Basis von $\elltwo$ und einer
      von $\Ltwo{\mathbb{R}}$.},
    voraussetzung={Abschnitt~\ref{c03.s7} bis Proposition~\ref{c03.prop.entwicklung};
      Vertiefungen~\ref{c03.v1} und~\ref{c03.v2}; die Vollständigkeit von $\elltwo$ und
      $\Ltwo{\mathbb{R}}$ aus Fundament~\ref{c03.f3}. Der Kern bis Kapitel~4 braucht den Block nicht.},
    zurueck={Auch in unendlicher Dimension hat jeder Hilbertraum eine Basis. Ist ihre
      Indexmenge abzählbar, so ist jeder Vektor die Summe seiner Komponenten, als Reihe, die in
      der Norm konvergiert. Im Kern genügt die endliche Summe.}]
  % Stimme: vl03.tex:212; v04 [26:54]-[27:16] („wir gehen davon aus, dass wir immer eine Basis
  % haben, das wird ab sofort angenommen in unserer Vorlesung, sonst ist das zu schwer“).  Ohne
  % den nachträglichen Rahmen „In der Vorlesung sage ich“ (docs/review/lectures/REVIEW-lec3.md,
  % Zeile VOICE (de.tex:1045)).  Abzählbare Indexmengen für die Summen: Reparatur nach
  % REVIEW-lec3.md, Zeile WRONG (de.tex:1048), und evidence/derivations/c03/REVIEW-lec3.md, B,
  % Zeile WRONG; Satz c03.thm.parseval und Proposition c03.prop.maximal sind auf abzählbare
  % Indexmengen beschränkt, D86 sagt, was eine Summe von Vektoren dann heißt.
  Wir gehen ab sofort davon aus, dass wir immer eine Basis haben, sonst wird es zu schwer. Für
  Hilberträume ist das mehr als eine Annahme, nämlich ein Satz. Zuerst muss man aber sagen,
  was eine Basis bei unendlich vielen Vektoren ist. Bei einer abzählbaren Indexmenge, etwa
  $X = \{1, 2, 3, \dots\}$, werden die Summen zu Reihen, und die Vollständigkeit aus
  \defref{D85} verlangt dann, dass Reihen konvergieren. Eine Indexmenge darf aber auch
  überabzählbar sein. Reed und Simon definieren die Basis deshalb ohne Summe:

  \tpcdef{D86}

  \begin{theorem}[Existenz einer Basis]\label{c03.thm.existenz}\claim[proof-omitted]{c03.cl.existenz}{source-quoted}
    Jeder Hilbertraum hat eine Orthonormalbasis.
  \end{theorem}

  Reed und Simon beweisen das mit dem Zornschen Lemma \srcref{reed-simon-1}{S.~44--45}. Der
  nächste Satz sagt, dass die Entwicklung nach einer Basis auch dann gilt, wenn sie
  unendlich viele Vektoren hat, und zwar abzählbar viele.

  \begin{theorem}[Entwicklung nach einer Orthonormalbasis]\label{c03.thm.parseval}\claim[proof-omitted]{c03.cl.parseval}{source-quoted}
    Sei $\{\ket{\phi_j}\}_{j \in X}$ eine Orthonormalbasis eines Hilbertraums $\Hil$ mit
    abzählbarer Indexmenge $X$. Dann gilt für jeden Vektor $\ket{\Psi} \in \Hil$
    \begin{equation}\label{c03.eq.parseval}
      \ket{\Psi} = \sum_{j \in X} \braket{\phi_j}{\Psi}\, \ket{\phi_j} ,
      \qquad
      \norm{\Psi}^2 = \sum_{j \in X} \abs{\braket{\phi_j}{\Psi}}^2 .
    \end{equation}
    Die erste Summe ist eine Reihe von Vektoren; sie konvergiert in der Norm, unabhängig von
    der Reihenfolge der Summanden, gegen $\ket{\Psi}$. Umgekehrt konvergiert
    $\sum_{j \in X} c_j \ket{\phi_j}$ für komplexe Zahlen $c_j$ mit
    $\sum_{j \in X} \abs{c_j}^2 < \infty$ in der Norm gegen einen Vektor in $\Hil$.
  \end{theorem}

  Das ist Satz~II.6 bei Reed und Simon, im Fall einer abzählbaren Indexmenge; die zweite
  Gleichung heißt dort Parsevalsche Gleichung \srcref{reed-simon-1}{S.~45--46}. Reed und
  Simon formulieren den Satz für jede Indexmenge. Wir brauchen nur abzählbare: Die Basen von
  $\elltwo$ und $\Ltwo{\mathbb{R}}$ weiter unten sind mit $k = 1, 2, 3, \dots$ und
  $n = 0, 1, 2, \dots$ abgezählt. Damit passt die Definition des Kerns zu der von Reed und
  Simon:

  \begin{proposition}[Zwei Definitionen, eine Basis]\label{c03.prop.maximal}\claim[proof-included]{c03.cl.maximal}{derived}
    Ein orthonormales System $\{\ket{\phi_j}\}_{j \in X}$ eines Hilbertraums $\Hil$ mit
    abzählbarer Indexmenge $X$ ist genau dann eine Orthonormalbasis im Sinne von
    \defref{D86}, wenn
    $\norm{\psi}^2 = \sum_{j \in X} \abs{\braket{\phi_j}{\psi}}^2$ für jeden Vektor
    $\ket{\psi} \in \Hil$ gilt. Für endliches $X$ ist das die Basis aus \defref{D85}.
  \end{proposition}
  \begin{proof}
    Ist das System eine Orthonormalbasis, so gilt die Gleichung nach
    Satz~\ref{c03.thm.parseval}. Ist es keine, so gibt es ein orthonormales System, das alle
    $\ket{\phi_j}$ und noch einen weiteren Vektor $\ket{\xi}$ enthält. Dann ist
    $\norm{\xi} = 1$ und $\braket{\phi_j}{\xi} = 0$ für jedes $j$; jede endliche Teilsumme
    von $\sum_j \abs{\braket{\phi_j}{\xi}}^2$ ist null, und die Gleichung für
    $\ket{\psi} = \ket{\xi}$ hieße $1 = 0$.
  \end{proof}

  \begin{example}[Eine Basis von $\elltwo$]\label{c03.ex.elltwo-basis}\claim{c03.cl.elltwo-basis}{derived}
    Sei $\vb{u}_k \in \elltwo$ der Vektor mit einer $1$ im $k$-ten Eintrag und sonst lauter
    Nullen, $k = 1, 2, 3, \dots$ Dann ist $\vb{u}_j\adj\vb{u}_k = \kron{j}{k}$, und für jedes
    $\vb{v} \in \elltwo$ ist $\vb{u}_k\adj\vb{v} = v_k$, also
    \begin{equation}\label{c03.eq.elltwo-basis}
      \sum_{k=1}^{\infty} \abs{\vb{u}_k\adj\vb{v}}^2 = \sum_{k=1}^{\infty} \abs{v_k}^2 = \vb{v}\adj\vb{v} .
    \end{equation}
    Weil $\elltwo$ nach Satz~\ref{c03.thm.vollstaendig} ein Hilbertraum ist, bilden die
    $\vb{u}_k$ nach Proposition~\ref{c03.prop.maximal} eine Orthonormalbasis von $\elltwo$.
    Die Vollständigkeitsrelation sagt hier: Die Teilsummen $\sum_{k \leq N} v_k \vb{u}_k$, also
    die ersten $N$ Einträge von $\vb{v}$ und danach lauter Nullen, konvergieren gegen
    $\vb{v}$. Das sieht man auch direkt: Der Abstand zum Quadrat ist
    $\sum_{k > N} \abs{v_k}^2$, und das ist der Rest einer konvergenten Reihe.
  \end{example}

  Für $\Ltwo{\mathbb{R}}$ gibt Satz~\ref{c03.thm.existenz} eine Basis, weil der Raum in der
  Fassung von Fundament~\ref{c03.f2} nach Satz~\ref{c03.thm.vollstaendig} ein Hilbertraum ist. Reed und
  Simon nennen auch eine konkrete.

  \begin{theorem}[Hermite-Funktionen]\label{c03.thm.hermite}\claim[proof-omitted]{c03.cl.hermite}{source-quoted}
    Die Hermite-Funktionen $\phi_n$, $n = 0, 1, 2, \dots$, bilden eine Orthonormalbasis von
    $\Ltwo{\mathbb{R}}$. Die erste ist $\phi_0(x) = \pi^{-1/4}\, e^{-x^2/2}$.
  \end{theorem}

  Reed und Simon nennen sie auch die Wellenfunktionen des harmonischen Oszillators und
  verweisen für den Beweis auf Übungsaufgaben \srcref{reed-simon-1}{S.~142}; sie kommen in
  Kapitel~7 wieder. Dass $\phi_0$ normiert ist, folgt aus Gleichung~\eqref{c03.eq.gauss}:
  \begin{equation}\label{c03.eq.phinull}\claim{c03.cl.phinull}{derived}
    \braket{\phi_0}{\phi_0} = \pi^{-1/2} \int_{-\infty}^{\infty} e^{-x^2}\, \dd x
    = \pi^{-1/2} \sqrt{\pi} = 1 .
  \end{equation}
\end{vertiefung}

\subsection{Das Gram-Schmidt-Verfahren}\label{c03.s7.1}

% Stimme: vl03.tex:212; v04 [27:22]-[29:22] („diese Vektoren sind beliebig, völlig beliebig,
% irgendeine Liste, die müssen keine Eigenschaften erfüllen“).
% Nicht übernommen: v04 [27:58]-[28:01] „dieses Orthonormierungsverfahren liefert uns immer
% eine Basis“.  Das stimmt nur unter einer Bedingung, die das Skript nennt (vl03.tex:221); der
% Text nennt sie in Teil 4 der Proposition.
Wie bildet man in der Praxis Basen? Mit dem sogenannten Orthonormierungsverfahren von Gram und
Schmidt. Man startet mit irgendeiner Liste von Vektoren $\ket{\Phi_1}, \ket{\Phi_2}, \dots$;
diese Vektoren sind völlig beliebig und müssen keine Eigenschaften erfüllen. Daraus bestimmt
man nacheinander paarweise orthogonale Vektoren $\ket{\Psi_n}$, mit einer einzigen Formel.

\begin{proposition}[Gram-Schmidt-Verfahren]\label{c03.prop.gram-schmidt}\claim[proof-included]{c03.cl.gram-schmidt}{derived}
  Seien $\ket{\Phi_1}, \dots, \ket{\Phi_N}$ Vektoren eines Hilbertraums $\Hil$, und seien
  $\ket{\Psi_1}, \dots, \ket{\Psi_N}$ nacheinander durch
  \begin{equation}\label{c03.eq.gram-schmidt}
    \ket{\Psi_n} = \ket{\Phi_n} - \sum_{k=1}^{n-1} \ket{\Psi_k}\, \frac{\braket{\Psi_k}{\Phi_n}}{\norm{\Psi_k}^2}
  \end{equation}
  bestimmt, wobei Summanden mit $\ket{\Psi_k} = 0$ ausgelassen werden. Dann gilt:
  \begin{enumerate}
    \item $\braket{\Psi_k}{\Psi_n} = 0$ für $k < n$.
    \item Für jedes $m \leq N$ ist jedes $\ket{\Phi_n}$ mit $n \leq m$ eine Linearkombination
      der $\ket{\Psi_k}$ mit $k \leq m$, und umgekehrt.
    \item Die Vektoren $\ket{\Psi_n} / \norm{\Psi_n}$ mit $\ket{\Psi_n} \neq 0$ bilden ein
      orthonormales System.
    \item Steht nur der Nullvektor auf allen $\ket{\Phi_n}$ senkrecht, so ist dieses System
      eine Basis von $\Hil$.
  \end{enumerate}
\end{proposition}
\begin{proof}
  Zu~1 und~2 durch Induktion über $n$; für $n = 1$ ist $\ket{\Psi_1} = \ket{\Phi_1}$. Seien
  $\ket{\Psi_1}, \dots, \ket{\Psi_{n-1}}$ paarweise orthogonal. Für $j < n$ mit
  $\ket{\Psi_j} \neq 0$ bleibt in
  \begin{equation}\label{c03.eq.gs-beweis}
    \braket{\Psi_j}{\Psi_n}
    = \braket{\Psi_j}{\Phi_n} - \sum_{k} \braket{\Psi_j}{\Psi_k}\, \frac{\braket{\Psi_k}{\Phi_n}}{\norm{\Psi_k}^2}
    = \braket{\Psi_j}{\Phi_n} - \braket{\Psi_j}{\Phi_n} = 0
  \end{equation}
  von der Summe nur der Summand $k = j$ übrig; für $\ket{\Psi_j} = 0$ ist das Skalarprodukt
  ohnehin null. Nach Gleichung~\eqref{c03.eq.gram-schmidt} ist $\ket{\Phi_n}$ gleich
  $\ket{\Psi_n}$ plus einer Linearkombination der $\ket{\Psi_k}$ mit $k < n$, und
  $\ket{\Psi_n}$ ist $\ket{\Phi_n}$ minus einer Linearkombination der $\ket{\Psi_k}$ mit
  $k < n$, die nach Induktionsannahme Linearkombinationen der $\ket{\Phi_k}$ mit $k < n$ sind.
  Zu~3: Nach den Regeln für Zahlen ist $\norm{z\psi} = \abs{z}\, \norm{\psi}$, also haben die
  geteilten Vektoren die Norm $1$, und orthogonal sind sie nach~1. Zu~4: Seien
  $\ket{\xi_1}, \dots, \ket{\xi_M}$ die Vektoren aus~3, sei $\ket{\psi} \in \Hil$ und
  $\ket{r} = \ket{\psi} - \sum_m \braket{\xi_m}{\psi}\, \ket{\xi_m}$. Wie im Beweis von
  Proposition~\ref{c03.prop.entwicklung} ist $\braket{\xi_m}{r} = 0$ für jedes $m$, also
  auch $\braket{\Psi_k}{r} = 0$ für jedes $k$ und nach~2 $\braket{\Phi_n}{r} = 0$ für jedes
  $n$. Nach der Voraussetzung ist $\ket{r} = 0$. Jeder Vektor lässt sich also nach den
  $\ket{\xi_m}$ entwickeln, und nach Proposition~\ref{c03.prop.entwicklung} ist das System
  vollständig.
\end{proof}

% Stimme: v04 [31:20]-[34:16] („ich werde das sehr oft machen, dass ich einen Vektor mal eine
% skalare Größe habe“; „in der Definition stand immer an der linken Seite“; „aber Sie müssen
% damit leben“; „n hat gar nichts mit der Dimension von H zu tun“; Division durch null);
% vl03.tex:216 mit Fußnote.  „nichtnegative reelle Zahl“ statt „die Norm auch“ (komplex):
% evidence/derivations/c03/REVIEW-lec3.md, B, Zeile MINOR (de.tex:1186); D83.
Alle Größen in der Formel kann man ausrechnen: Das Skalarprodukt ist eine komplexe Zahl, die
Norm eine nichtnegative reelle Zahl, die $\ket{\Psi_k}$ sind Vektoren, und die Summe ist eine Summe von Vektoren. Die Zahl
steht hier rechts vom Vektor; in der Definition stand sie immer links. Das mache ich sehr oft,
und Sie müssen damit leben: Gemeint ist $\ket{\Psi_k}\, z = z \ket{\Psi_k}$. Der Index $n$ hat
gar nichts mit der Dimension des Raumes zu tun; er zählt nur die Liste ab. Die Summanden mit
$\ket{\Psi_k} = 0$ müssen wir weglassen, weil sonst $\norm{\Psi_k}^2 = 0$ im Nenner stünde und
wir durch null teilen würden. Es kann durchaus passieren, dass manche $\ket{\Psi_n}$ null sind
(Übung~\ref{c03.exc.gram-schmidt}).

% Stimme: vl03.tex:223-250; v04 [38:11]-[46:17] („mein Lieblings-Hilbertraum“; „diese beiden
% Vektoren habe ich gerade ausgedacht“).
Weil das Verfahren aus mehreren Teilen besteht und verwirrend sein kann, rechnen wir es einmal
ganz durch.

\begin{example}[Gram-Schmidt in $\mathbb{C}^2$]\label{c03.ex.gram-schmidt}\claim{c03.cl.gs-beispiel}{derived}
  Wir arbeiten in meinem Lieblings-Hilbertraum $\mathbb{C}^2$ und starten mit zwei Vektoren, die
  ich mir gerade ausgedacht habe:
  \begin{equation}\label{c03.eq.gs-start}
    \ket{\Phi_1} = \begin{pmatrix} 1 \\ 0 \end{pmatrix},
    \qquad
    \ket{\Phi_2} = \begin{pmatrix} 1 \\ 1 \end{pmatrix} .
  \end{equation}
  Im ersten Schritt läuft die Summe von $1$ bis $0$; es wird also gar nicht summiert, und
  $\ket{\Psi_1} = \ket{\Phi_1}$. Im zweiten Schritt bleibt nur der Summand $k = 1$. Mit
  $\braket{\Psi_1}{\Phi_2} = 1$ und $\norm{\Psi_1}^2 = 1$ ist
  \begin{equation}\label{c03.eq.gs-beispiel}
    \ket{\Psi_2} = \ket{\Phi_2} - \ket{\Psi_1}\, \frac{\braket{\Psi_1}{\Phi_2}}{\norm{\Psi_1}^2}
    = \begin{pmatrix} 1 \\ 1 \end{pmatrix} - \begin{pmatrix} 1 \\ 0 \end{pmatrix}
    = \begin{pmatrix} 0 \\ 1 \end{pmatrix} .
  \end{equation}
  Beide Vektoren haben schon die Norm eins: Aus der ausgedachten Liste ist die Basis
  $\ket{e_0}, \ket{e_1}$ geworden. Die Bedingung von Teil~4 der Proposition ist erfüllt: Aus
  $\braket{\Phi_1}{\psi} = \psi_1 = 0$ und $\braket{\Phi_2}{\psi} = \psi_1 + \psi_2 = 0$ folgt
  $\ket{\psi} = 0$.
\end{example}

% Stimme: vl03.tex:250; v04 [46:17]-[46:45] („im dreidimensionalen Fall geht es genauso, es wird
% nur aufwendiger“).
In höheren Dimensionen geht die Rechnung genauso, sie wird nur aufwendiger und aufwendiger.
Das Verfahren funktioniert immer, in dem Sinn, dass es immer ein orthonormales System liefert;
eine Basis ist es, wenn die Bedingung von Teil~4 erfüllt ist.

\begin{exercise}[Ein Vektor fällt weg]\label{c03.exc.gram-schmidt}\claim{c03.cl.gs-null}{derived}
  Nehmen Sie zu den beiden Vektoren aus Beispiel~\ref{c03.ex.gram-schmidt} den dritten Vektor
  $\ket{\Phi_3} = \begin{pmatrix} 1 \\ -1 \end{pmatrix}$ hinzu und berechnen Sie $\ket{\Psi_3}$.
\end{exercise}
\begin{solution}
  Mit $\ket{\Psi_1} = \ket{e_0}$, $\ket{\Psi_2} = \ket{e_1}$, $\braket{\Psi_1}{\Phi_3} = 1$
  und $\braket{\Psi_2}{\Phi_3} = -1$ ist
  \begin{equation}\label{c03.eq.gs-null}
    \ket{\Psi_3} = \begin{pmatrix} 1 \\ -1 \end{pmatrix} - \begin{pmatrix} 1 \\ 0 \end{pmatrix}
    - (-1) \begin{pmatrix} 0 \\ 1 \end{pmatrix} = \begin{pmatrix} 0 \\ 0 \end{pmatrix} .
  \end{equation}
  Der dritte Vektor ist schon eine Linearkombination der ersten beiden,
  $\ket{\Phi_3} = 2 \ket{\Phi_1} - \ket{\Phi_2}$, und fällt weg; das orthonormale System bleibt
  $\ket{e_0}, \ket{e_1}$.
\end{solution}

% ── Fundament ─────────────────────────────────────────────────────────────────────────
% Neue Abschnitte des Kapitels (vl04, vl05) kommen VOR diese Zeile.

\begin{fundament}[id=c03.f1,
    titel={Die Cauchy-Schwarz-Ungleichung},
    wozu={Vertiefung~\ref{c03.v1} zeigt mit einer groben Abschätzung, dass das
      Skalarprodukt auf $\elltwo$ endlich ist. Hier steht die scharfe Schranke: Der Betrag
      des Skalarprodukts ist höchstens das Produkt der Längen.},
    voraussetzung={Proposition~\ref{c03.prop.skalarprodukt}, die Schwarz'sche Ungleichung
      (Proposition~\ref{c03.prop.schwarz}) und Vertiefung~\ref{c03.v1}. Kern und
      Vertiefungen brauchen den Block nicht.},
    zurueck={Für zwei Vektoren ist $\abs{\vb{v}\adj\vb{w}}^2$ nie größer als
      $(\vb{v}\adj\vb{v})(\vb{w}\adj\vb{w})$, in $\mathbb{C}^n$ wie in $\elltwo$.}]
  % Der Beweis in Spalten, der den von c03.prop.schwarz wiederholte, ist durch die
  % Spezialisierung ersetzt (docs/review/lectures/REVIEW-lec3.md, Zeile MINOR (de.tex:1255)).
  \begin{lemma}[Cauchy-Schwarz in $\mathbb{C}^n$]\label{c03.lem.cs}\claim[proof-linked]{c03.cl.cs}{derived}
    Für alle $\vb{v}, \vb{w} \in \mathbb{C}^n$ gilt
    \begin{equation}\label{c03.eq.cs}
      \abs{\vb{v}\adj\vb{w}}^2 \leq (\vb{v}\adj\vb{v})\,(\vb{w}\adj\vb{w}) .
    \end{equation}
  \end{lemma}

  Das Lemma ist die Schwarz'sche Ungleichung aus Proposition~\ref{c03.prop.schwarz} für das
  Skalarprodukt $\vb{v}\adj\vb{w}$ von $\mathbb{C}^n$. Dieses erfüllt die Bedingungen von
  \defref{D82}, denn das sind die Regeln~2 bis~4 von Proposition~\ref{c03.prop.skalarprodukt};
  mit $\norm{v}^2 = \vb{v}\adj\vb{v}$ ist Gleichung~\eqref{c03.eq.cs} das Quadrat von
  Gleichung~\eqref{c03.eq.schwarz}. Reed und Simon beweisen dieselbe Ungleichung für jeden
  Vektorraum mit Skalarprodukt, auf anderem Weg \srcref{reed-simon-1}{S.~38}. Für $\elltwo$
  folgt sie hier aus dem Lemma durch einen Grenzübergang.

  \begin{proposition}[Cauchy-Schwarz in $\elltwo$]\label{c03.prop.cs-elltwo}\claim[proof-included]{c03.cl.cs-elltwo}{derived}
    Für alle $\vb{v}, \vb{w} \in \elltwo$ gilt
    $\abs{\vb{v}\adj\vb{w}}^2 \leq (\vb{v}\adj\vb{v})\,(\vb{w}\adj\vb{w})$.
  \end{proposition}
  \begin{proof}
    Für jedes $N$ ist Lemma~\ref{c03.lem.cs}, angewandt auf die ersten $N$ Einträge, die
    Ungleichung
    $\abs{\sum_{j \leq N} v_j^* w_j}^2 \leq \sum_{j \leq N} \abs{v_j}^2 \sum_{j \leq N} \abs{w_j}^2$.
    Die rechte Seite ist höchstens $(\vb{v}\adj\vb{v})(\vb{w}\adj\vb{w})$. Die linke Seite
    strebt für $N \to \infty$ gegen $\abs{\vb{v}\adj\vb{w}}^2$, weil die Reihe nach
    Proposition~\ref{c03.prop.elltwo} konvergiert.
  \end{proof}
\end{fundament}

\begin{fundament}[id=c03.f2,
    titel={Funktionen, die sich nur auf einer Nullmenge unterscheiden},
    wozu={Vertiefung~\ref{c03.v2} lässt offen, ob $\braket{f}{f} = 0$ nur für $f = 0$ gilt.
      Für Funktionen stimmt das nicht. Hier steht die Reparatur und was sie für den Wert
      einer Funktion an einer einzelnen Stelle bedeutet.},
    voraussetzung={Vertiefung~\ref{c03.v2}. Lebesgue-Maß und Lebesgue-Integral werden nur
      dem Namen nach benutzt.},
    zurueck={Zwei Funktionen, die sich nur an einzelnen Stellen unterscheiden, sind
      derselbe Vektor in $\Ltwo{\mathbb{R}}$. Der Wert an einer einzelnen Stelle hat keine
      Bedeutung, wohl aber Integrale über Intervalle.}]
  \begin{example}[Eine Funktion mit Länge null]\label{c03.ex.punkt}\claim{c03.cl.punkt}{derived}
    Sei $f(0) = 1$ und $f(x) = 0$ für $x \neq 0$. Diese Funktion ist stückweise stetig und
    nicht die Nullfunktion, aber
    \begin{equation}\label{c03.eq.punkt}
      \braket{f}{f} = \int_{-\infty}^{\infty} \abs{f(x)}^2\, \dd x = 0 ,
    \end{equation}
    denn ein Integrand, der nur an einer Stelle von null verschieden ist, hat das Integral
    null. Regel~4 von Proposition~\ref{c03.prop.skalarprodukt} gilt für Funktionen also
    nicht.
  \end{example}

  Die Reparatur besteht darin, solche Funktionen mit der Nullfunktion zu identifizieren. In
  der genauen Fassung ist $\Ltwo{\mathbb{R}}$ deshalb ein Raum von Klassen von Funktionen.

  \tpcdef{D80}

  Für quadratintegrierbare $f$ und $g$ gilt
  \begin{equation}\label{c03.eq.klassen}\claim{c03.cl.klassen}{derived}
    f \sim g
    \quad \Leftrightarrow \quad
    \int_{-\infty}^{\infty} \abs{f(x) - g(x)}^2\, \dd x = 0 .
  \end{equation}
  Reed und Simon sagen für integrierbare Funktionen, dass $f(x) = g(x)$ fast überall dasselbe
  ist wie $\int \abs{f - g}\, \dd x = 0$ \srcref{reed-simon-1}{S.~17}. Angewandt auf die
  integrierbare Funktion $\abs{f - g}^2$ und die Nullfunktion gibt das
  Gleichung~\eqref{c03.eq.klassen}, denn $\abs{f - g}^2$ ist genau dort null, wo $f = g$ ist.
  Daraus folgt auch, dass das Skalarprodukt nicht davon abhängt, welche Funktion man aus
  einer Klasse nimmt: Ist $n \sim 0$, so ist $\abs{n^* g}$ fast überall null, nach derselben
  Aussage also $\int \abs{n^* g}\, \dd x = 0$ und damit $\braket{n}{g} = 0$. In $\Ltwo{\mathbb{R}}$ gilt nun auch Regel~4: Aus
  $\braket{f}{f} = 0$ folgt $f \sim 0$, und das ist der Nullvektor.

  Der Preis ist, dass ein Element von $\Ltwo{\mathbb{R}}$ keinen Wert an einer einzelnen
  Stelle hat: Ändert man eine Funktion an einer Stelle, so bleibt sie in ihrer Klasse.
  Reed und Simon schreiben das für die integrierbaren Funktionen so: Für ein Element $f$ hat
  das Symbol $f(x)$ an einer bestimmten Stelle $x$ keinen Sinn \srcref{reed-simon-1}{S.~17}.
  Für Funktionen, die einen stetigen Vertreter haben, hilft die folgende Aussage:

  \begin{proposition}[Stetige Vertreter]\label{c03.prop.stetig}\claim[proof-omitted]{c03.cl.stetig}{source-quoted}
    Enthält eine Klasse eine stetige Funktion, so enthält sie keine andere stetige
    Funktion. Man darf dann $f(x)$ für den Wert dieser stetigen Funktion schreiben.
  \end{proposition}

  Reed und Simon sagen das für die Klassen integrierbarer Funktionen
  \srcref{reed-simon-1}{S.~18}; es hängt nur an der Relation „fast überall gleich“, die für
  $\Ltwo{\mathbb{R}}$ dieselbe ist.
\end{fundament}

\begin{fundament}[id=c03.f3,
    titel={Vollständigkeit},
    wozu={Der Kern sagt, dass einem Vektorraum mit Skalarprodukt zum Hilbertraum noch eine
      Bedingung fehlt und dass sie in $\mathbb{C}^n$ immer gilt. Hier steht, was die
      Bedingung verlangt, warum $\mathbb{C}^n$, $\elltwo$ und $\Ltwo{\mathbb{R}}$ sie
      erfüllen, und ein Raum, der sie verletzt.},
    voraussetzung={Vertiefungen~\ref{c03.v1} und~\ref{c03.v2}, Fundament~\ref{c03.f2};
      Folgen und Grenzwerte.},
    zurueck={$\mathbb{C}^n$, $\elltwo$ und $\Ltwo{\mathbb{R}}$ sind vollständig. Wer mit
      ihnen rechnet, braucht über die Bedingung nicht nachzudenken.}]
  \tpcdef{D79}

  Das Vorbild ist der Unterschied zwischen den rationalen und den reellen Zahlen: Eine Folge
  rationaler Zahlen, die sich einer irrationalen Zahl nähert, ist eine Cauchy-Folge, hat in
  den rationalen Zahlen aber keinen Grenzwert. $\mathbb{R}$ ist vollständig, $\mathbb{Q}$
  nicht \srcref{reed-simon-1}{S.~6}.

  \begin{proposition}[$\mathbb{C}^n$ ist vollständig]\label{c03.prop.cn-vollstaendig}\claim[proof-included]{c03.cl.cn-vollstaendig}{derived}
    Jede Cauchy-Folge in $\mathbb{C}^n$ konvergiert gegen einen Vektor in $\mathbb{C}^n$.
  \end{proposition}
  \begin{proof}
    Sei $\vb{v}^{(1)}, \vb{v}^{(2)}, \dots$ eine Cauchy-Folge. Für jeden Index $j$ ist
    \begin{equation}\label{c03.eq.komponente}
      \abs{v_j^{(m)} - v_j^{(k)}}^2
      \leq \sum_{l=1}^{n} \abs{v_l^{(m)} - v_l^{(k)}}^2
      = d\bigl(\vb{v}^{(m)}, \vb{v}^{(k)}\bigr)^2 .
    \end{equation}
    Also sind die Real- und die Imaginärteile der $j$-ten Einträge Cauchy-Folgen reeller
    Zahlen; weil $\mathbb{R}$ vollständig ist, haben sie Grenzwerte. Sei $v_j$ der
    zusammengesetzte Grenzwert und $\vb{v}$ der Vektor mit diesen Einträgen. Dann ist
    $d(\vb{v}^{(m)}, \vb{v})^2 = \sum_{j=1}^{n} \abs{v_j^{(m)} - v_j}^2$ eine Summe von
    $n$ Gliedern, von denen jedes gegen null geht.
  \end{proof}

  Dass die Bedingung nicht leer ist, zeigt ein Raum, der zwischen $\mathbb{C}^n$ und
  $\elltwo$ liegt.

  \begin{example}[Abbrechende Folgen sind nicht vollständig]\label{c03.ex.endlich}\claim{c03.cl.endlich}{derived}
    Sei $F$ die Menge der Vektoren mit unendlich vielen Einträgen, von denen nur endlich
    viele von null verschieden sind, mit dem Skalarprodukt von $\elltwo$. Mit $F$ liegen
    auch Summen und Vielfache in $F$. Sei $\vb{v}^{(N)}$ der Vektor mit den Einträgen
    $1, \tfrac12, \dots, \tfrac1N$ und danach lauter Nullen. Für $M > N$ ist nach
    Gleichung~\eqref{c03.eq.rate}
    \begin{equation}\label{c03.eq.endlich}
      d\bigl(\vb{v}^{(M)}, \vb{v}^{(N)}\bigr)^2
      = \sum_{j=N+1}^{M} \frac{1}{j^2}
      \leq \sum_{j=N+1}^{M} \Bigl( \frac{1}{j-1} - \frac{1}{j} \Bigr)
      = \frac{1}{N} - \frac{1}{M} < \frac{1}{N} ,
    \end{equation}
    die Folge ist also eine Cauchy-Folge. Hätte sie einen Grenzwert $\vb{u}$ in $F$, so wäre
    nach dem Argument von Gleichung~\eqref{c03.eq.komponente} für jedes $j$
    $u_j = \lim_{N} v_j^{(N)} = 1/j$, und $\vb{u}$ hätte unendlich viele Einträge ungleich
    null. Also konvergiert die Folge in $F$ nicht. In $\elltwo$ dagegen konvergiert sie,
    gegen den Vektor mit $v_j = 1/j$ aus Beispiel~\ref{c03.ex.rate}: Dieselbe Rechnung
    mit $M \to \infty$ gibt den Abstand höchstens $1/\sqrt{N}$.
  \end{example}

  Der Raum $\elltwo$ füllt die Lücken von $F$ also gerade auf; Reed und Simon sagen
  allgemeiner, dass die abbrechenden Folgen in diesen Räumen dicht liegen
  \srcref{reed-simon-1}{S.~69}.

  \begin{theorem}[Vollständigkeit von $\elltwo$ und $\Ltwo{\mathbb{R}}$]\label{c03.thm.vollstaendig}\claim[proof-omitted]{c03.cl.vollstaendig}{source-quoted}
    Der Raum $\elltwo$ ist vollständig. Der Raum $\Ltwo{\mathbb{R}}$ in der Fassung von
    \defref{D80} ist vollständig (Satz von Riesz und Fischer).
  \end{theorem}

  Die Quelle ist Reed und Simon: für $\elltwo$ als der Fall $p = 2$ der Folgenräume, die
  sie für jedes $p \geq 1$ betrachten, für $\Ltwo{\mathbb{R}}$ als der Fall $p = 2$ des
  Satzes von Riesz und Fischer für die Räume $L^p$ \srcref{reed-simon-1}{S.~68--69}. Im Fall von $\Ltwo{\mathbb{R}}$ gilt der
  Satz für die Fassung aus Fundament~\ref{c03.f2}, mit messbaren Funktionen und dem
  Lebesgue-Integral.

  % Stimme: v03 [87:13]-[89:12].
  Der Versuchung, Grenzfälle zu betrachten, werden wir widerstehen. Ich werde in dieser
  Vorlesung nur legale Rechnungen durchführen, zumindest hoffe ich das. Physiker gehen mit
  diesen Regeln ohnehin ein bisschen lockerer um als Mathematiker: Sobald etwas unendlich
  wird, sagen wir, wir haben einen Fehler gemacht. Unendlich darf keine Vorhersage sein,
  deshalb lassen wir diese Rechnung weg. Wenn Sie aber ein bisschen in die mathematische
  Richtung gehen möchten, dann sind diese Bedingungen wirklich sehr wichtig.

  Mathematisch ist das keine Regel. Ein Vektor kann nämlich in
  $\elltwo$ liegen und trotzdem eine unendliche Summe liefern: Für $v_j = 1/j$ ist
  $\sum_j \abs{v_j}^2$ endlich, aber $\sum_j j\, \abs{v_j}^2 = \sum_j 1/j$ divergiert
  (Beispiel~\ref{c03.ex.rate}). Summen dieser Art treten später bei Erwartungswerten auf.
  Ob eine solche Rechnung erlaubt ist, muss man dann prüfen; das Werkzeug dafür sind die
  Definitionsbereiche von Operatoren.
\end{fundament}
